\documentclass[11pt]{article}
\usepackage{tcee}
\addbibresource{3A41.bib}
\tcetema[La inflación: causas, efectos, hiperinflación y deflación]{3.A.41}{La inflación: causas y efectos sobre la eficiencia económica y el bienestar. Hiperinflación y deflación.}
\tcefecha{05/11/2026}
\begin{document}
\tceetitulo

\section*{Introducción}
\minutos{2,5}

\begin{itemize}
\item \textbf{\emph{Enganche}}:
  \begin{itemize}
  \item «Quien juega con fuego, se quema», dice el refrán. En octubre de 1923, en Alemania, los precios se multiplicaron casi por 300 en un solo mes: se duplicaban cada 3,7 días\nocite{HankeKrus2013World}. El Gobierno pagaba su déficit con la imprenta, y el marco acabó sin valor\nocite{Sargent1982Ends}.
  \item Un siglo después, la preocupación de la zona del euro era la contraria: los precios bajaron en tasa interanual en 2009, entre diciembre de 2014 y marzo de 2015, en 2016 y de agosto a diciembre de 2020, y el BCE luchaba contra la deflación\nocite{Eurostat2026IAPC}.
  \item Y el fuego ha vuelto: en octubre de 2026 la inflación de la zona del euro ha subido al 4,0~\%, empujada por la energía, que se encarece un 20,8~\% en un año\nocite{Eurostat2026IAPCoct}; en España se mantiene en el 4,9~\%\nocite{INE2026IPCAdelantadoOct}. Las expectativas a largo plazo, sin embargo, siguen ancladas en el 2,0~\%\nocite{BCE2026SPF4T}.
  \item Detrás hay 2 preguntas: \emph{¿de dónde sale ese fuego y cuánto calor conviene?}
  \end{itemize}

\item \textbf{\emph{Relevancia}}:
  \begin{itemize}
  \item \emph{Para la teoría}: la inflación es el campo de pruebas de la macroeconomía. La neutralidad del dinero, la curva de \autor{Phillips}, las expectativas racionales, la inconsistencia temporal, las rigideces de precios de la nueva economía keynesiana y la relación entre la política monetaria y la fiscal se discutieron, sobre todo, a propósito de ella.
  \item \emph{Para la política económica}: la estabilidad de precios es el objetivo principal del Sistema Europeo de Bancos Centrales\nocite{TFUE2016}, que el BCE concreta en un objetivo simétrico del 2~\% a medio plazo, fijado en julio de 2021 y confirmado en junio de 2025\nocite{BCE2021Estrategia,BCE2025Estrategia} \vertema{3.A.36}.
  \item \emph{Para el bienestar}: la inflación es un impuesto que no aprueba ningún parlamento, y recae más sobre los hogares de renta baja, que gastan una parte mayor de su renta en energía y alimentos\nocite{BCE2022RentaBaja}.
  \end{itemize}

\item \textbf{\emph{Contextualización}}:
  \begin{itemize}
  \item \emph{Definiciones}:
    \begin{itemize}
    \item la \concepto{inflación} es el aumento \emph{sostenido} y \emph{generalizado} del nivel de precios; se mide con la tasa de variación de un índice de precios, $\pi_t\equiv(P_t-P_{t-1})/P_{t-1}$;
    \item la \concepto{deflación} es su contrario, una caída sostenida y generalizada (no basta con un mes negativo), y la \concepto{desinflación}, que la inflación baje sin dejar de ser positiva;
    \item la \concepto{hiperinflación} es, en la definición de \textcite{Cagan1956Monetary}, una inflación de más del 50~\% \emph{al mes}.
    \end{itemize}
  \item \emph{Medición}: el IPC (y en la Unión Europea el IAPC, que permite comparar países) mide los precios de consumo; el deflactor del PIB, los de todo lo que se produce en el país, sin lo importado; y la \concepto{inflación subyacente} quita los componentes más volátiles, la energía y los alimentos, para ver la tendencia \vertema{3.A.9} \vertema{4.A.1}. Todos tienen sesgos: el IPC tiende a sobrestimar la subida del coste de la vida, porque no recoge del todo las mejoras de calidad, los productos nuevos ni la sustitución entre bienes (\ref{sec:inflacion-optima}).
  \item \emph{Las 4 etapas de la teoría}:
    \begin{enumerate}
    \item la \emph{teoría cuantitativa}: a largo plazo, la cantidad de dinero fija los precios \vertema{3.A.3};
    \item la \emph{keynesiana y la síntesis}: inflación de demanda y de costes, y la curva de \autor{Phillips} como menú entre inflación y paro\nocite{Phillips1958,SamuelsonSolow1960Analytical} \vertema{3.A.5};
    \item la de las \emph{expectativas}, de 1968 a los 80: la curva aumentada\nocite{Friedman1968Role,Phelps1967Phillips}, las expectativas racionales\nocite{Lucas1972Expectations}, la inconsistencia temporal\nocite{KydlandPrescott1977Rules} y, en el lado fiscal, la aritmética monetarista desagradable\nocite{SargentWallace1981};
    \item la de la \emph{nueva síntesis}: precios rígidos con expectativas racionales y bancos centrales independientes con un objetivo de inflación. Entre 1999 y 2020 la inflación media de la zona del euro fue del 1,6~\%\nocite{Eurostat2026IAPC}; tras 2008 la preocupación fue la deflación, y desde 2021, la vuelta de la inflación\nocite{BernankeBlanchard2025What}.
    \end{enumerate}
  \item \emph{Consenso y debate}:
    \begin{itemize}
    \item hay \emph{consenso} en 2 cosas: a largo plazo la inflación sigue al dinero, y una inflación alta y volátil es mala para la eficiencia, el bienestar y la estabilidad;
    \item el \emph{debate} está en qué la mueve a corto plazo y por qué persiste (para los modelos neoclásicos, la inflación desaparece al retirar el estímulo; para los keynesianos, las rigideces la mantienen), y en qué objetivo conviene: la regla de \autor{Friedman} pide una deflación suave; las rigideces nominales, inflación nula; el límite inferior de los tipos y la rigidez de los salarios a la baja, algo positiva.
    \end{itemize}
  \end{itemize}

\item \textbf{\emph{Problemática}}:
  \begin{itemize}
  \item ¿Qué causa la inflación: el dinero, la demanda, los costes o su propia inercia? ¿Es siempre un fenómeno monetario?
  \item ¿Qué cuesta la inflación, depende de que se anticipe y existe una tasa óptima?
  \item ¿Por qué son tan dañinos los extremos, la hiperinflación y la deflación, y cómo se sale de ellos?
  \end{itemize}

\item \textbf{\emph{Estructura}}:
  \begin{itemize}
  \item La exposición sigue el orden del título, en 3 partes:
  \begin{itemize}
  \item \concepto{I. Causas de la inflación} (\ref{sec:causas}): del debate entre clásicos y keynesianos al triángulo de \autor{Gordon}, con sus 3 vértices, la demanda (de la política monetaria, de la fiscal y de otras perturbaciones), los costes y la inercia, y su aplicación a la inflación de 2021-2026;
  \item \concepto{II. Efectos sobre la eficiencia económica y el bienestar} (\ref{sec:efectos}): los costes de la inflación anticipada y de la no anticipada, la tasa de inflación óptima y la estabilización;
  \item \concepto{III. Contextos extremos: hiperinflación y deflación} (\ref{sec:extremos}): los 2 incendios, el de demasiado fuego y el de ninguno.
  \end{itemize}
  Los extremos van al final porque se entienden con las causas (el déficit que se paga con dinero, la demanda que se hunde) y con los costes, que en ellos se multiplican.
  \end{itemize}
\end{itemize}

\begin{anotaciones}
Para seguir la inflación: el IPC del INE, \url{https://www.ine.es/dyngs/INEbase/es/operacion.htm?c=Estadistica_C&cid=1254736176802}, y el IAPC de Eurostat, \url{https://ec.europa.eu/eurostat/web/hicp}. Sobre el origen de la inflación de 2021: Banco de España, «El impacto de las perturbaciones de oferta y de demanda sobre la evolución reciente de la actividad y los precios», \emph{Informe trimestral de la economía española}, 4/2021, recuadro 3\nocite{BdE2021Recuadro3}. Para ver: los debates de la Fundación Ramón Areces «El retorno de la inflación», \url{https://youtu.be/M8aQHlDLhIM}, y «La inflación y la crisis del coste de la vida», \url{https://youtu.be/899qcwdx_Gc}. Para leer: la historia del monetarismo, con frases de sus autores, en \href{https://www.hetwebsite.net/het/essays/monetarism/monetarism.htm}{\emph{The History of Economic Thought Website}, «Monetarism»}.
\end{anotaciones}

\begin{esquema}[unbreakable]
\textbf{3.A.41. La inflación: causas, efectos, hiperinflación y deflación}\\[2pt]
Introducción\esqtiempo{2,5}\\
I. Causas de la inflación\esqtiempo{13}\\
\hspace*{1.5em}1. Debate tradicional de la inflación en modelos clásicos y keynesianos\\
\hspace*{1.5em}2. Causas de la inflación en modelos dinámicos: el triángulo de \autor{Gordon} (1982)\\
\hspace*{3em}política monetaria · política fiscal · NEK: demanda y costes · inercia\\
\hspace*{1.5em}3. La inflación de 2021-2026 a la luz del triángulo\\
II. Efectos sobre la eficiencia económica y el bienestar\esqtiempo{8}\\
\hspace*{1.5em}1. Costes de la inflación anticipada y no anticipada\\
\hspace*{1.5em}2. ¿Existe una tasa de inflación óptima?\\
\hspace*{1.5em}3. Inflación y estabilización macroeconómica\\
III. Contextos extremos: hiperinflación y deflación\esqtiempo{5,5}\\
\hspace*{1.5em}1. Hiperinflación\\
\hspace*{1.5em}2. Deflación\\
Conclusión\esqtiempo{1}
\end{esquema}

\begin{notaopositor}
Reparto del cante: introducción, \tiempo{2,5}; I, \tiempo{13} (I.1, \tiempo{1,5}; I.2, \tiempo{9,5}: presentación del triángulo, \tiempo{1}; política monetaria, \tiempo{3}; política fiscal, \tiempo{1,5}; las demás perturbaciones de demanda y las de costes en la NEK, \tiempo{2}; inercia, \tiempo{2}; I.3, \tiempo{2}); II, \tiempo{8} (II.1, \tiempo{4}; II.2, \tiempo{3}; II.3, \tiempo{1}); III, \tiempo{5,5} (III.1, \tiempo{3,5}; III.2, \tiempo{2}), y conclusión, \tiempo{1}. Lo que la exposición resume (los modelos clásico y keynesiano y la historia de la curva de \autor{Phillips}, el desarrollo de \autor{Lucas}, de \autor{Barro} y \autor{Gordon} y de la aritmética monetarista, las rigideces y la persistencia en detalle, la inflación óptima con impuestos y con precios rígidos, \autor{Cagan} con expectativas racionales y la inflación y las cuentas públicas) está en los anexos~\ref{anexo:clasico-keynesiano} a~\ref{anexo:cuentas-publicas}; las cajas grises son para consultar.
\end{notaopositor}

% ==========================================================================
\section{Causas de la inflación}\label{sec:causas}
\minutos{13}

\begin{itemize}
\item Empezamos por las causas, del debate tradicional entre clásicos y keynesianos a los modelos dinámicos.
\item \emph{Conclusión anticipada}: a largo plazo, la inflación es un fenómeno monetario; a corto plazo la explican 3 fuerzas, la demanda, los costes y la inercia, y la política monetaria y la fiscal deciden si el dinero la alimenta y, por tanto, si se queda.
\end{itemize}

% --------------------------------------------------------------------------
\subsection{Debate tradicional de la inflación en modelos clásicos y keynesianos}\label{sec:debate-tradicional}

\begin{itemize}
\item Partimos de la pregunta más antigua: ¿de dónde sale una subida \emph{sostenida} de los precios? La respuesta tradicional separa el largo plazo, en el que manda el dinero, del corto plazo, en el que importa si los precios y los salarios se ajustan.
\item \emph{De la identidad a la teoría cuantitativa}: la \concepto{ecuación de cambios}, $M\,V\equiv P\,Y$, es una \emph{identidad} contable del gasto: se cumple siempre, porque la velocidad $V$ se define como el cociente entre el gasto nominal y el dinero, y no dice nada de causas. La \concepto{teoría cuantitativa del dinero} de \autor{Irving Fisher} (1911) es esa misma identidad reinterpretada con supuestos causales: la velocidad es estable (la fijan los hábitos de pago y las instituciones), la producción está en su nivel natural (la fijan factores reales), el dinero es exógeno y la causalidad va del dinero a los precios \vertema{3.A.3} \vertema{3.A.35}. Con esos supuestos, la ecuación en tasas,
  \begin{equation}\label{eq:cuantitativa}
  \gamma_M+\gamma_V=\pi+\gamma_Y ,
  \end{equation}
  dice que la inflación sigue al crecimiento del dinero menos el de la producción. Para \textcite{Friedman1963Inflation}, «la inflación es siempre y en todas partes un fenómeno monetario», en el sentido de que solo puede producirla un aumento de la cantidad de dinero más rápido que el de la producción\nocite{Friedman1970Counter} \vertema{3.A.5}.\footnote{\autor{Milton Friedman} recibió el Premio Nobel en 1976 «por sus logros en el análisis del consumo, la historia y la teoría monetarias y por demostrar la complejidad de la política de estabilización»\nocite{Nobel1976Friedman}.}
\item \emph{Evidencia}: a largo plazo se cumple bien. Con datos de 110 países entre 1960 y 1990, la correlación entre el crecimiento del dinero y la inflación es de 0,925 con la base monetaria, 0,958 con M1 y 0,950 con M2, y la del dinero con el crecimiento real, casi nula\nocite{McCandlessWeber1995} \vertema{3.A.35}.
\item \emph{El debate tradicional, en una viñeta}\nocite{Serletis2007Demand}:
  \begin{itemize}
  \item para los \emph{clásicos}, con precios y salarios flexibles, el dinero cumple 2 propiedades distintas:
    \begin{itemize}
    \item la \concepto{neutralidad}: cualquier variación de la cantidad de dinero cambia en la misma proporción todas las variables nominales (precios, salarios nominales) sin afectar a las reales (producción, empleo, tipo de interés real, salario real);
    \item la \concepto{dicotomía clásica}: las variables reales se determinan en un bloque de ecuaciones en el que no entra el dinero (el mercado de trabajo, la función de producción, el ahorro y la inversión), y el mercado de dinero solo fija el nivel de precios: «el dinero es como un velo», que cubre las variables reales sin tocarlas.
    \end{itemize}
    Un estímulo monetario solo sube los precios, y la inflación desaparece al retirarlo;
  \item para los \emph{keynesianos}, con el salario nominal rígido a la baja, se rompen las 2: la demanda mueve primero la producción y el empleo, y la inflación puede venir tanto de la demanda como de los costes (la puja entre salarios y beneficios) \vertema{3.A.4};
  \item para los \emph{monetaristas}, la inflación sostenida es siempre monetaria, y el corto plazo se explica por las expectativas \vertema{3.A.5}.
  \end{itemize}
\end{itemize}

\begin{anotaciones}
\textbf{Neutralidad y dicotomía: por qué son independientes} (para consultar; en el cante, una frase):
\begin{itemize}
\item \emph{Son preguntas distintas}:
  \begin{itemize}
  \item la \emph{neutralidad} es un resultado de \emph{estática comparativa}: compara 2 equilibrios que solo se diferencian en la cantidad de dinero y pregunta qué ha cambiado. Es neutral si el nuevo equilibrio es el viejo con todas las variables nominales multiplicadas por el mismo número;
  \item la \emph{dicotomía} es una propiedad de la \emph{estructura} del modelo: si el sistema de ecuaciones se puede partir en 2 bloques y el real se resuelve solo, sin saber nada del dinero ni de los precios.
  \end{itemize}
\item \emph{Por qué en el modelo clásico se dan las 2}: en las ecuaciones las variables nominales solo aparecen divididas por otra nominal ($W/P$, $M/P$), así que multiplicar $M$, $P$ y $W$ por el mismo número no cambia nada (neutralidad), y además las ecuaciones del mercado de trabajo, la producción y el ahorro e inversión no contienen $M/P$ (dicotomía).
\item \emph{Neutralidad sin dicotomía}, el caso importante. Si el consumo depende de la riqueza real, y en ella entran los saldos reales, $C=C(Y,r,M/P)$ (el \emph{efecto saldos reales} de \textcite{Patinkin1956Money}), el bloque real ya no se puede resolver sin el mercado de dinero: no hay dicotomía. Pero si $M$ se duplica y $P$ se duplica, $M/P$ no cambia y el equilibrio real tampoco: el dinero sigue siendo neutral. Ejemplo numérico: con $M=100$ y $P=1$ los saldos reales son 100; con $M=200$, el único equilibrio es $P=2$, los saldos reales vuelven a 100 y el consumo, la producción y el tipo real son los de antes. Lo mismo pasa en el modelo con dinero en la utilidad (MIU) de \textcite{Sidrauski1967Rational}: los saldos reales están en la utilidad (no hay dicotomía) y, sin embargo, el dinero es neutral \vertema{3.A.35}. \autor{Patinkin} mostró que la dicotomía «pura» de los clásicos era incluso incoherente, porque dejaba el nivel de precios indeterminado; lo que se sostiene es la neutralidad.
\item \emph{Dicotomía sin neutralidad}: si el bloque real no depende del dinero, pero la demanda de dinero tiene ilusión monetaria (por ejemplo, el público demanda una cantidad \emph{nominal} que no crece en proporción a los precios), un aumento de $M$ no cambia las variables reales, pero el nivel de precios y el salario nominal suben en otra proporción que el dinero: hay dicotomía y no hay neutralidad tal como se ha definido.
\item \emph{Resumen}: la dicotomía con una demanda de dinero homogénea (en saldos reales) basta para la neutralidad, pero la neutralidad no exige dicotomía. Y la \emph{superneutralidad} es otra cosa: que el \emph{ritmo} de crecimiento del dinero tampoco afecte a las variables reales (\ref{sec:anticipada}).
\end{itemize}
\end{anotaciones}

\begin{itemize}
\item En los 3 enfoques, un impulso que se da una vez (más dinero, más gasto, una caída de la productividad) sube el \emph{nivel} de precios una vez. Para que haya inflación, el impulso tiene que repetirse o algo tiene que perpetuarlo: es lo que estudian los modelos dinámicos. Las ecuaciones de los modelos clásico y keynesiano están en el anexo~\ref{anexo:clasico-keynesiano}.
\end{itemize}

% --------------------------------------------------------------------------
\subsection{Causas de la inflación en modelos dinámicos:\cortetitulo el triángulo de \autor{Gordon} (1982)}\label{sec:triangulo-gordon}

\begin{itemize}
\item Una vez visto que a largo plazo la inflación sigue al dinero, pasamos a lo que la mueve de un año para otro. El instrumento es la curva de \autor{Phillips}, que para \textcite{Akerlof2002Behavioral} es «probablemente la relación macroeconómica más importante».\footnote{Lo dijo en su conferencia del Nobel, en vídeo en \href{https://www.nobelprize.org/prizes/economic-sciences/2001/akerlof/lecture/}{nobelprize.org} (minuto 20:32). \autor{George Akerlof} compartió el Premio Nobel de 2001 con \autor{Michael Spence} y \autor{Joseph Stiglitz} «por sus análisis de los mercados con información asimétrica»\nocite{Nobel2001AkerlofSpenceStiglitz}.} Con expectativas, la \concepto{curva aumentada} es \vertema{3.A.28}
  \begin{equation}\label{eq:phillips-aumentada}
  \pi_t=\pi^e_t-a\,\big(U_t-U^n\big)+\varepsilon_t,\qquad a>0,
  \end{equation}
  con $U^n$ la tasa natural de paro (o NAIRU), $a$ la pendiente y $\varepsilon_t$ las perturbaciones de oferta. A largo plazo, con $\pi=\pi^e$ y sin perturbaciones, es vertical en $U^n$.
\item \textcite{Gordon1982Price} ordena las causas de la inflación a partir de esta curva en el \concepto{modelo del triángulo}\nocite{Gordon1977Can,Gordon1997Time,Gordon2011Phillips}: la inflación tiene 3 orígenes, \emph{demanda}, \emph{costes} e \emph{inercia}, y cada uno es un término de la ecuación (gráfico~\ref{fig:esquema-inflacion}). Si el público forma sus expectativas con la inflación de los $k$ periodos anteriores, cada uno con su peso,
  \funcionclave[eq:triangulo-gordon]{\pi_t=\underbrace{\sum_{i=1}^{k}a_i\,\pi_{t-i}}_{\text{inercia}}\;\underbrace{-\,a\,\big(U_t-U^n\big)}_{\text{demanda}}\;+\underbrace{\varepsilon_t}_{\text{costes}},\qquad\sum_{i=1}^{k}a_i=1}
  con $a_i$ los pesos de la inflación pasada (locales: las $a_i$ con subíndice no son la pendiente $a$). La suma empieza en el periodo anterior, $i=1$: la inflación de hoy está a la izquierda.
  \begin{itemize}
  \item \concepto{Inflación de demanda} (\emph{demand-pull}): el paro por debajo de la tasa natural, la producción por encima de la potencial, empuja los precios hacia arriba.
  \item \concepto{Inflación de costes} (\emph{cost-push}): una perturbación de oferta (energía, materias primas, márgenes, tipo de cambio) sube los precios con cualquier nivel de paro.
  \item \concepto{Inflación inercial} (\emph{built-in}): la inflación que ya estaba en marcha sigue por las expectativas, la indexación y los contratos.
  \end{itemize}
\item Como los pesos suman 1, el triángulo es aceleracionista: con el paro en su tasa natural y sin perturbaciones, la inflación se queda donde estaba, y solo baja si el paro supera a $U^n$. Por eso desinflar cuesta (\ref{sec:desinflacion}). Con histéresis, cuenta también la variación del paro \vertema{3.A.28}.
\end{itemize}

\begin{anotaciones}
\textbf{De la curva aumentada al triángulo, paso a paso} (para consultar; no se canta):
\begin{itemize}
\item \emph{Paso 1, la curva aumentada}:
  \begin{itemize}
  \item se parte de~\eqref{eq:phillips-aumentada}: la inflación es la esperada, menos el efecto del paro sobre la tasa natural, más las perturbaciones de oferta.
  \end{itemize}
\item \emph{Paso 2, las expectativas miran atrás}:
  \begin{itemize}
  \item la inflación esperada es una media ponderada de la pasada, $\pi^e_t=\sum_{i\geq1}a_i\,\pi_{t-i}$, con pesos que suman 1;
  \item las adaptativas son un caso particular: con $\pi^e_t=\pi^e_{t-1}+\gamma(\pi_{t-1}-\pi^e_{t-1})$ y sustituyendo hacia atrás, $a_i=\gamma(1-\gamma)^{i-1}$, que suman 1.
  \end{itemize}
\item \emph{Paso 3, se sustituye}:
  \begin{itemize}
  \item llevando las expectativas a la curva aumentada sale~\eqref{eq:triangulo-gordon}, con un término de cada vértice.
  \end{itemize}
\item \emph{Paso 4, la forma general con retardos}: \autor{Gordon} deja que también la demanda y la oferta actúen con retraso y escribe, con el operador de retardos $L$ ($L\,\pi_t=\pi_{t-1}$),
  \begin{equation*}
  \pi_t=\underbrace{a(L)\,\pi_{t-1}}_{\text{inercia}}+\underbrace{b(L)\,D_t}_{\text{demanda}}+\underbrace{c(L)\,z_t}_{\text{costes}}+e_t ,
  \end{equation*}
  \begin{itemize}
  \item con $D_t$ una medida del exceso de demanda (la brecha de paro o de producción), $z_t$ las perturbaciones de oferta medidas (precios relativos de la energía y de las importaciones, productividad, controles de precios) y $e_t$ un error; $a(L)$, $b(L)$ y $c(L)$ son polinomios de retardos (locales; $z_t$ y $e_t$, también);
  \item en sus estimaciones trimestrales la inercia entra con muchos retardos, varios años de inflación pasada.
  \end{itemize}
\item \emph{Paso 5, la restricción de la tasa natural}:
  \begin{itemize}
  \item si los pesos de la inflación pasada suman 1, $a(1)=1$, en un estado estacionario sin perturbaciones $b(1)\,D=0$: el paro solo puede estar en $U^n$, y la curva de largo plazo es vertical;
  \item si sumaran menos de 1, habría un menú permanente, $\pi=b(1)\,D/[1-a(1)]$, que es lo que \autor{Friedman} y \autor{Phelps} negaron.
  \end{itemize}
\end{itemize}
\end{anotaciones}

\begin{figure}[H]
\centering
\caption{De dónde viene la inflación: el triángulo de \autor{Gordon}}\label{fig:esquema-inflacion}
\grafico[0.55\textwidth]{esquema-origen-inflacion}
\fuente{Elaboración propia a partir de \textcite{Gordon1982Price,Gordon2011Phillips}. Los 3 vértices son los 3 orígenes de la inflación y los 3 términos de~\eqref{eq:triangulo-gordon}.}
\end{figure}

\begin{itemize}
\item Cada vértice es una causa y pide una respuesta distinta. Empezamos por la demanda que crean las 2 políticas, la monetaria (\ref{sec:politica-monetaria}) y la fiscal (\ref{sec:politica-fiscal}); después, con el modelo canónico de la nueva economía keynesiana, las otras perturbaciones de demanda y las de costes (\ref{sec:nek-origenes}) y la inercia (\ref{sec:inercia}). Es el esquema de la exposición y la forma en que los bancos centrales leen los datos.
\end{itemize}

% ..........................................................................
\subsubsection{Inflación de demanda: la política monetaria}\label{sec:politica-monetaria}

\paragraph{De la curva de \autor{Phillips} a la hipótesis aceleracionista}

\begin{itemize}
\item \emph{La aportación original}: \textcite{Phillips1958} encontró, con datos del Reino Unido de 1861 a 1957, una relación inversa y no lineal entre el crecimiento de los salarios nominales y el paro: con poco paro, los salarios subían deprisa.\footnote{No fue el primero: \textcite{Fisher1926Statistical} ya había relacionado la inflación con el empleo en Estados Unidos. Su artículo se volvió a publicar en 1973 con el título «I Discovered the Phillips Curve»\nocite{Fisher1973Discovered}.} Era una regularidad empírica, sin teoría detrás \vertema{3.A.28}.
\item \emph{La síntesis}: \textcite{SamuelsonSolow1960Analytical} la pasaron a precios, con un margen constante sobre los costes laborales, y la leyeron como un \emph{menú} de política económica: más inflación a cambio de menos paro, de forma permanente \vertema{3.A.5}.\footnote{\autor{Paul Samuelson} recibió el Premio Nobel en 1970 «por el trabajo científico con el que ha desarrollado la teoría económica estática y dinámica y ha contribuido activamente a elevar el nivel de análisis de la ciencia económica»\nocite{Nobel1970Samuelson}, y \autor{Robert Solow}, en 1987, «por sus contribuciones a la teoría del crecimiento económico»\nocite{Nobel1987Solow}.} La historia completa está en el anexo~\ref{anexo:phillips-historia}.
\item \emph{La crítica monetarista}: \textcite{Friedman1968Role} y \textcite{Phelps1967Phillips} la corrigieron: los trabajadores negocian salarios reales, así que lo que cuenta es la inflación \emph{esperada}, y sale la curva aumentada~\eqref{eq:phillips-aumentada}.\footnote{\autor{Edmund Phelps} recibió el Premio Nobel en 2006 «por su análisis de las disyuntivas intertemporales en la política macroeconómica»\nocite{Nobel2006Phelps}.} Una expansión monetaria (o cualquier aumento de la demanda) baja el paro por debajo de $U^n$ y sube la inflación por encima de la esperada (gráfico~\ref{fig:phillips-aumentada}): la economía pasa de $A$ a $B$ por la curva de corto plazo. Es echar leña al fuego: estimular la demanda cuando la economía ya está en su potencial solo aviva la inflación.
\item \emph{Con expectativas adaptativas},
  \begin{equation*}
  \pi^e_t=\pi^e_{t-1}+\gamma\,\big(\pi_{t-1}-\pi^e_{t-1}\big),\qquad 0<\gamma\leq1,
  \end{equation*}
  el público acaba esperando la inflación más alta y la curva de corto plazo sube (de $B$ a $C$). Con $\gamma=1$, $\pi^e_t=\pi_{t-1}$, y la curva aumentada es la \concepto{curva aceleracionista} \vertema{3.A.28}:
  \begin{equation}\label{eq:aceleracionista}
  \Delta\pi_t\equiv\pi_t-\pi_{t-1}=-a\,\big(U_t-U^n\big).
  \end{equation}
  Mantener el paro por debajo de la tasa natural exige una inflación cada año mayor (de $C$ a $D$): la demanda que se repite se convierte en inercia, y el menú de \autor{Samuelson} y \autor{Solow} solo existe a corto plazo.
\item Por eso la inflación sostenida de demanda es, a la larga, monetaria: sin un dinero que crezca, la demanda no puede seguir empujando año tras año~\eqref{eq:cuantitativa}.
\end{itemize}

\begin{figure}[H]
\centering
\caption{La curva de \autor{Phillips} aumentada con expectativas: corto y largo plazo}\label{fig:phillips-aumentada}
\grafico[0.6\textwidth]{phillips-aumentada}
\fuente{Elaboración propia a partir de \textcite{Friedman1968Role} y \textcite{Phelps1967Phillips}. Una curva de corto plazo por cada inflación esperada (la primera es el menú de \textcite{SamuelsonSolow1960Analytical}); la de largo plazo es vertical en $U^n$.}
\end{figure}

\paragraph{De las islas de \autor{Phelps} a la oferta de \autor{Lucas}}

\subparagraph{a) Idea}

\begin{itemize}
\item Las adaptativas permiten que el público se equivoque año tras año en la misma dirección. \textcite{Lucas1972Expectations} pone expectativas racionales y explica por qué el dinero mueve la producción a corto plazo aunque nadie se equivoque de forma sistemática: por la \emph{información imperfecta} \vertema{3.A.6}.\footnote{\autor{Robert Lucas} recibió el Premio Nobel en 1995 «por haber desarrollado y aplicado la hipótesis de las expectativas racionales, transformando así el análisis macroeconómico y profundizando en la comprensión de la política económica»\nocite{Nobel1995Lucas}.}
\end{itemize}

\subparagraph{b) Modelo}

\hspace*{3.6em}\emph{i) Supuestos}

\begin{itemize}
\item Los productores viven en mercados separados, las \concepto{islas} de \textcite{Phelps1970Microeconomic}: cada uno ve el precio de su isla, $p_{i,t}=p_t+z_{i,t}$, con $z_{i,t}$ una perturbación real de la isla, de media cero, pero no el nivel general $p_t$ (en logaritmos).
\item Solo quieren producir más si sube su precio \emph{relativo}; las expectativas son racionales.
\end{itemize}

\hspace*{3.6em}\emph{ii) Desarrollo}

\begin{itemize}
\item Cuando su precio sube, cada productor no sabe si es una perturbación real (le conviene producir más) o una subida general (no le conviene). Atribuye a causas reales la parte $\omega=\sigma^2_z/(\sigma^2_z+\sigma^2_p)$ de la subida, el \emph{coeficiente de proyección}: la parte de la varianza de su precio que se debe a las perturbaciones reales.
\item Agregando las islas, la \concepto{oferta de \autor{Lucas}}\nocite{Lucas1973Some}:
  \begin{equation}\label{eq:oferta-lucas}
  y_t=y^n+\psi\,\big(p_t-p^e_t\big),\qquad\psi\equiv\omega\,\varepsilon_S ,
  \end{equation}
  con $y^n$ la producción natural y $\varepsilon_S$ la elasticidad de la oferta de cada isla a su precio relativo. Solo las sorpresas de precios mueven la producción. El desarrollo, en el anexo~\ref{anexo:lucas}.
\end{itemize}

\hspace*{3.6em}\emph{iii) Implicaciones}

\begin{itemize}
\item La política monetaria puede tener efectos reales, pero solo con 3 condiciones:
  \begin{enumerate}
  \item \emph{a muy corto plazo}: en cuanto los productores ven el nivel general, corrigen; el efecto no persiste;
  \item \emph{por sorpresa}: la política sistemática se incorpora a $p^e_t$ y no mueve la producción, la \concepto{proposición de inefectividad}\nocite{SargentWallace1975Rational};
  \item \emph{con poca variabilidad nominal}: donde la demanda nominal es más volátil, $\omega$ es menor, la curva es casi vertical y la inflación compra menos producción\nocite{Lucas1973Some}.\footnote{Con precios rígidos sale algo parecido por otra vía: la curva es más vertical cuanto más alta es la inflación media, porque los precios se revisan más a menudo\nocite{BallMankiwRomer1988} \vertema{3.A.7}.}
  \end{enumerate}
\item Como $\psi$ cambia con la política, la curva no sirve como un menú estable, la \concepto{crítica de \autor{Lucas}} \vertema{3.A.6}. Y con inflación errática los precios dejan de informar: es uno de los costes de la inflación sobre la eficiencia (\ref{sec:costes-inflacion}).
\end{itemize}

\subparagraph{c) Valoración}

\begin{itemize}
\item Da fundamento a la curva de \autor{Phillips} con expectativas racionales, pero no explica por qué los efectos del dinero duran años si las estadísticas de precios salen cada mes: es la puerta a las rigideces de la nueva economía keynesiana (\ref{sec:nek-origenes}).
\end{itemize}

\paragraph{Reglas frente a discrecionalidad: \autor{Kydland} y \autor{Prescott} (1977) y \autor{Barro} y \autor{Gordon} (1983)}

\subparagraph{a) Idea}

\begin{itemize}
\item Si solo la inflación por sorpresa baja el paro y el público lo sabe, ¿por qué ha habido tanta inflación? \textcite{KydlandPrescott1977Rules} y \textcite{BarroGordon1983Positive} responden con una causa \emph{institucional}: la \concepto{inconsistencia temporal} de la política monetaria \vertema{3.A.6} \vertema{3.A.36}.\footnote{\autor{Finn Kydland} y \autor{Edward Prescott} recibieron el Premio Nobel en 2004 «por sus contribuciones a la macroeconomía dinámica: la consistencia temporal de la política económica y las fuerzas que impulsan los ciclos económicos»\nocite{Nobel2004KydlandPrescott}.}
\end{itemize}

\subparagraph{b) Modelo}

\hspace*{3.6em}\emph{i) Supuestos}

\begin{itemize}
\item La curva de \autor{Phillips} con sorpresas, $U_t=U^n-\alpha\,(\pi_t-\pi^e_t)$, con $\alpha=1/a$.
\item Las \emph{preferencias del banco central}: minimiza una pérdida cuadrática en el paro y la inflación, y quiere un paro $U^*$ \emph{por debajo} del natural (porque el paro natural es ineficiente, por los impuestos o el poder de mercado):
  \begin{equation*}
  L_t=(U_t-U^*)^2+\gamma\,(\pi_t-\pi^*)^2,\qquad U^*<U^n,
  \end{equation*}
  con $\gamma$ el peso de la inflación (aquí $\gamma$ no es la de las adaptativas).
\item Un juego en 2 etapas: primero el público fija $\pi^e_t$ y después el banco central elige $\pi_t$ con discrecionalidad.
\end{itemize}

\hspace*{3.6em}\emph{ii) Desarrollo}

\begin{itemize}
\item Fijadas las expectativas, sorprender con inflación baja el paro hacia $U^*$, y el banco central tiene la tentación de hacerlo. El público lo sabe y espera la inflación con la que el banco central ya no quiere sorprender: en equilibrio el paro queda en $U^n$ y la inflación, por encima del objetivo, en la \concepto{inflación de discrecionalidad}
  \funcionclave[eq:sesgo]{\pi^{\mathrm D}=\pi^*+\frac{\alpha}{\gamma}\,\big(U^n-U^*\big)>\pi^*}
  El sesgo inflacionista es mayor cuanto más quiere el banco central bajar el paro ($U^n-U^*$), cuanto más responde el paro a las sorpresas ($\alpha$) y cuanto menos pesa la inflación en su pérdida ($\gamma$). La derivación y los 3 puntos (regla, tentación y discrecionalidad), en el anexo~\ref{anexo:barro-gordon}.
\end{itemize}

\hspace*{3.6em}\emph{iii) Implicaciones}

\begin{itemize}
\item La inflación no compra ningún paro: es el precio de la falta de compromiso. La regla ($\pi^*$ con $U^n$) sería mejor que la discrecionalidad, pero no es creíble.
\item \emph{Qué se puede hacer}:
  \begin{itemize}
  \item \emph{reglas} con compromiso legal;
  \item \emph{reputación} en el juego repetido \vertema{3.A.14};
  \item un banco central \emph{independiente y conservador}, con un $\gamma$ mayor que el de la sociedad\nocite{Rogoff1985Optimal};
  \item \emph{atarse las manos} con un tipo de cambio fijo frente a un país de baja inflación\nocite{GiavazziPagano1988Advantage};
  \item \emph{objetivos de inflación}, que empezaron en Nueva Zelanda en 1990\nocite{BernankeEtAl1997NuevaZelanda} \vertema{3.A.36}.
  \end{itemize}
  La misma inconsistencia temporal aparece en la regulación \vertema{3.A.20}.
\end{itemize}

\subparagraph{c) Valoración}

\begin{itemize}
\item Explica la inflación de los 70 y la ola de independencia de los bancos centrales y de objetivos de inflación de los 90. Su límite: los bancos centrales independientes no parecen querer un paro por debajo del natural, y hoy el problema ha sido a menudo el contrario, una inflación por debajo del objetivo.
\end{itemize}

% ..........................................................................
\subsubsection{Inflación de demanda: la política fiscal}\label{sec:politica-fiscal}

\begin{itemize}
\item \emph{La restricción presupuestaria une las 2 políticas}: la inflación es monetaria, pero el dinero puede estar al servicio del déficit. Lo que no cubren los impuestos se paga con deuda o con dinero, y emitir dinero da al Estado un ingreso real, el \concepto{señoreaje}. Si el Estado tiene poco acceso a la deuda o un sistema fiscal pobre, el banco central acaba financiándolo; cuánto puede recaudar así y qué pasa si se le pide demasiado es la hiperinflación (\ref{sec:hiperinflacion}) \vertema{3.A.39}.
\end{itemize}

\begin{anotaciones}
\textbf{La restricción presupuestaria consolidada del sector público} (para consultar; en el cante, la última ecuación y su lectura)\nocite{Walsh2017Monetary} \vertema{3.A.39}:
\begin{itemize}
\item \emph{El Tesoro}: $G_t+i_{t-1}D^T_{t-1}=T_t+(D^T_t-D^T_{t-1})+\mathrm{RBC}_t$, con $G_t$ el gasto sin intereses, $D^T$ la deuda total, $T_t$ los impuestos y $\mathrm{RBC}_t$ los beneficios que le transfiere el banco central.
\item \emph{El banco central}: $(D^M_t-D^M_{t-1})+\mathrm{RBC}_t=i_{t-1}D^M_{t-1}+(M_t-M_{t-1})$: compra deuda ($D^M$) y paga beneficios con los intereses que cobra y con la base monetaria nueva.
\item \emph{Consolidando}, con $D\equiv D^T-D^M$ la deuda en manos del público:
  \begin{equation*}
  G_t+i_{t-1}D_{t-1}=T_t+(D_t-D_{t-1})+(M_t-M_{t-1}).
  \end{equation*}
\item \emph{En porcentaje del PIB}, dividiendo por $P_tY_t$, con $g$, $\tau$, $d$ y el señoreaje $s_t\equiv(M_t-M_{t-1})/(P_tY_t)$, y el tipo real y el crecimiento real:
  \begin{equation*}
  d_t-d_{t-1}=\underbrace{(g_t-\tau_t)}_{\text{déficit primario}}+\underbrace{\frac{r_{t-1}-\gamma_Y}{1+\gamma_Y}\,d_{t-1}}_{\text{bola de nieve}}-\underbrace{s_t}_{\text{señoreaje}}-\underbrace{\frac{(\pi_t-\mathbb{E}_{t-1}[\pi_t])\,(1+r_{t-1})}{(1+\pi_t)(1+\gamma_Y)}\,d_{t-1}}_{\text{inflación no esperada}}.
  \end{equation*}
\item \emph{Lectura}: la deuda sobre el PIB sube con el déficit primario y sola si el tipo real supera al crecimiento; baja con el señoreaje y con la inflación no esperada, que reduce el valor real de la deuda nominal. La inflación aparece 2 veces: como ingreso (señoreaje) y como impuesto sobre los acreedores (sorpresa).
\end{itemize}
\end{anotaciones}

\begin{itemize}
\item De la restricción salen 3 formas en que la política fiscal causa inflación:
  \begin{enumerate}
  \item \emph{La monetización directa}: el banco central compra la deuda o da crédito al Tesoro y el dinero crece con el déficit. Es la vía de las hiperinflaciones (\ref{sec:hiperinflacion}); por eso el artículo~123 del TFUE prohíbe financiar a los Estados\nocite{TFUE2016}.
  \item \concepto{La aritmética monetarista desagradable}\nocite{SargentWallace1981}:\footnote{\autor{Thomas Sargent} compartió el Premio Nobel de 2011 con \autor{Christopher Sims} «por su investigación empírica sobre la causa y el efecto en la macroeconomía»\nocite{Nobel2011SargentSims}.}
    \begin{itemize}
    \item \emph{supuesto}: la política fiscal es la que manda (\concepto{dominación fiscal}): fija una senda de déficits que no se ajusta, y el público no compra deuda más allá de un límite;
    \item \emph{el dilema del banco central}: o monetiza hoy, con algo de inflación hoy, o se niega y el déficit se financia con deuda, que crece con la bola de nieve si $r>\gamma_Y$, hasta el límite; a partir de ahí tiene que monetizar una deuda mayor, con más señoreaje y más inflación que si lo hubiera hecho desde el principio;
    \item \emph{con expectativas racionales}: el público anticipa la inflación futura, demanda menos dinero hoy y la inflación sube \emph{ya}: apretar hoy la política monetaria puede dar más inflación incluso hoy, y el banco central no controla la inflación ni siquiera a largo plazo \vertema{3.A.39};
    \item \emph{lectura}: quién cede es un \emph{juego de la gallina} entre el Tesoro y el banco central\nocite{BarthelemyPlantin2019Fiscal} \vertema{3.A.14}. La derivación, en el anexo~\ref{anexo:fiscal}.
    \end{itemize}
  \item \concepto{La teoría fiscal del nivel de precios}\nocite{Leeper1991,Sims1994Simple,Woodford1995Price}:
    \begin{itemize}
    \item lleva la idea más lejos: no hace falta que el dinero crezca. Si los superávits primarios futuros ($\tau-g$, sobre el PIB) están dados y no se ajustan para pagar la deuda, con una política fiscal activa y una monetaria pasiva en la terminología de \textcite{Leeper1991}, lo que se ajusta es el nivel de precios, que reduce el valor real de la deuda nominal hasta igualarlo al valor actual de los superávits,
      \begin{equation}\label{eq:valoracion}
      \frac{(1+i_{t-1})\,D_{t-1}}{P_t\,Y_t}=\sum_{k=0}^{\infty}\left(\frac{1+\gamma_Y}{1+r}\right)^{k}\left(\tau_{t+k}-g_{t+k}+s_{t+k}\right),
      \end{equation}
      con $D_{t-1}$ la deuda nominal heredada, $s$ el señoreaje sobre el PIB, $r$ el tipo real y $\gamma_Y$ el crecimiento \vertema{3.A.39};
    \item una expansión fiscal que nadie espera pagar con impuestos sube los precios aunque el dinero no crezca; \textcite{Cochrane2023Fiscal} lee así la inflación de Estados Unidos de 2021-2022;\footnote{Hay otra rama, minoritaria, en la que las variables fiscales afectan al nivel de precios a través de la demanda de dinero.}
    \item la crítica principal es que trata como condición de equilibrio lo que es una restricción presupuestaria del Gobierno\nocite{Buiter2002Fiscal}.
    \end{itemize}
  \end{enumerate}
\item La lección común: un banco central independiente necesita una política fiscal sostenible.
\end{itemize}

\begin{anotaciones}
Una introducción a la teoría fiscal del nivel de precios de uno de sus autores: \autor{John H. Cochrane}, «Fiscal Theory of the Price Level», conferencia en el Robert Schuman Centre del Instituto Universitario Europeo, \url{https://youtu.be/nfilAvVx4WI?t=175}.
\end{anotaciones}

% ..........................................................................
\subsubsection{Las demás perturbaciones de demanda y las de costes en el modelo canónico de la NEK}\label{sec:nek-origenes}

\begin{itemize}
\item Vistas las 2 políticas, el resto de los orígenes de la inflación se explican con el \concepto{modelo canónico de la nueva economía keynesiana} (NEK): empresas con poder de mercado que solo revisan su precio de vez en cuando (con probabilidad $\eta$ de \emph{no} revisarlo cada periodo)\nocite{Calvo1983Staggered,Gali2015Monetary}, hogares que optimizan y expectativas racionales \vertema{3.A.7}. Cómo modela la NEK estas rigideces (\autor{Fischer}, \autor{Taylor} y \autor{Calvo}) está en el anexo~\ref{anexo:rigideces}. Son 3 ecuaciones log-linealizadas alrededor de un estado estacionario sin inflación:
  \begin{itemize}
  \item la \emph{IS dinámica}, la demanda: $x_t=\mathbb{E}_t[x_{t+1}]-\frac{1}{\theta}\big(i_t-\mathbb{E}_t[\pi_{t+1}]-r^n_t\big)$;
  \item la \emph{curva de \autor{Phillips} neokeynesiana} (NKPC), la oferta, con la perturbación de costes:
    \funcionclave[eq:nkpc]{\pi_t=\beta\,\mathbb{E}_t[\pi_{t+1}]+\kappa\,x_t+\underbrace{u_t}_{\text{costes}}}
  \item la \emph{regla de \autor{Taylor}}, el banco central: $i_t=\rho+\phi_\pi\,\pi_t+\phi_x\,x_t+v_t$;
  \end{itemize}
  con $x_t$ la brecha de producción, $\theta$ la inversa de la elasticidad intertemporal del consumo, $r^n_t$ el tipo de interés natural, $\beta$ el factor de descuento, $\kappa$ la pendiente, $\rho$ la tasa de descuento, $\phi_\pi$ y $\phi_x$ las respuestas del banco central y $v_t$ una perturbación monetaria. La curva de \autor{Phillips} sale de la optimización de las empresas (anexo~\ref{anexo:nkpc}).
\item \emph{Los 3 orígenes, en las 3 ecuaciones}: la demanda entra por la IS ($r^n_t$) y la política monetaria por la regla ($v_t$); los costes, por la NKPC ($u_t$); y la inercia, por lo que la NKPC \emph{no} tiene, la inflación pasada (\ref{sec:inercia}).
\end{itemize}

\paragraph{Otras perturbaciones de demanda}

\begin{itemize}
\item La demanda no solo sube con el dinero o el déficit: un aumento de la confianza o de la riqueza de los hogares, el gasto del ahorro acumulado, un auge de la inversión, una mayor demanda exterior o una depreciación hacen lo mismo.
\item En el modelo todas suben el tipo de interés natural, $r^n_t$: si el banco central no sube el tipo nominal lo suficiente, la brecha de producción se abre y, por la NKPC, la inflación sube.
\item Con una perturbación de demanda, la producción y la inflación suben a la vez, y el banco central puede estabilizar las 2 con el tipo de interés: es la \concepto{divina coincidencia} (gráfico~\ref{fig:nek-xpi}, panel a) \vertema{3.A.7}.
\end{itemize}

\paragraph{Perturbaciones de costes (\emph{cost-push})}

\begin{itemize}
\item Una vez vista la inflación que tira de la demanda, pasamos a la que empuja desde los costes: la producción baja y los precios suben a la vez, la \concepto{estanflación}. El término de oferta se generalizó en las curvas de \autor{Phillips} con las crisis del petróleo de los 70\nocite{Gordon2011Phillips}.
\item En la NKPC, $u_t$ es la \concepto{perturbación de costes}: todo lo que sube el coste marginal o el margen deseado sin que cambie la brecha (la energía, las materias primas, los impuestos indirectos, los márgenes). Es la misma perturbación que $\varepsilon_t$ en~\eqref{eq:phillips-aumentada} y~\eqref{eq:triangulo-gordon}, con la letra de cada modelo.
\item También un modelo de ciclo real (RBC), con precios y salarios flexibles, recoge una perturbación de oferta: una caída de la productividad baja la producción y sube el nivel de precios. Pero allí la inflación no tiene un papel propio, porque lo que importa son las variables reales, y el banco central no tiene ningún dilema \vertema{3.A.42}. Por eso la inflación de costes se estudia con rigideces nominales.
\item \emph{El dilema}: con $u_t>0$ la NKPC sube, y el banco central no puede estabilizar a la vez la inflación y la actividad (gráfico~\ref{fig:nek-xpi}, panel b) \vertema{3.A.36}. Si la perturbación es transitoria y las expectativas están ancladas, puede «mirar a través» de ella; si amenaza con efectos de segunda ronda, tiene que subir los tipos aunque la actividad se resienta.
\end{itemize}

\begin{figure}[H]
\centering
\caption{Perturbaciones de demanda y de costes en la nueva economía keynesiana}\label{fig:nek-xpi}
\grafico[0.85\textwidth]{nek-xpi}
\fuente{Elaboración propia a partir de \textcite[cap.~5]{Gali2015Monetary}. a) De demanda: la inflación y la brecha se mueven juntas y el tipo de interés las estabiliza a la vez. b) De costes: el dilema.}
\end{figure}

\begin{anotaciones}
\textbf{Por qué el modelo canónico se puede dibujar en el plano $(x_t,\pi_t)$, y con qué supuestos} (para consultar; no se canta):
\begin{itemize}
\item \emph{En sentido estricto, no se puede}: el modelo canónico es un sistema de ecuaciones en diferencias estocásticas con expectativas, no un par de curvas estáticas. La NKPC y la IS contienen $\mathbb{E}_t[\pi_{t+1}]$ y $\mathbb{E}_t[x_{t+1}]$, que son endógenas, y no hay una «curva de oferta» ni una «de demanda» que se desplacen.
\item \emph{Los supuestos que hay detrás del gráfico}:
  \begin{itemize}
  \item las ecuaciones están \emph{linealizadas} alrededor de un \emph{estado estacionario} sin inflación ni brecha, y solo valen para desviaciones pequeñas;
  \item las perturbaciones siguen procesos \emph{AR(1)}, por ejemplo $u_t=\rho_u\,u_{t-1}+\epsilon^u_t$, con $0\leq\rho_u<1$;
  \item se busca la solución de \emph{variables de estado mínimas}, de la forma $x_t=\psi_x\,u_t$ y $\pi_t=\psi_\pi\,u_t$ (con $\psi_x$ y $\psi_\pi$ por determinar);
  \item de ahí, \emph{expectativas proporcionales}: $\mathbb{E}_t[x_{t+1}]=\rho_u\,x_t$ y $\mathbb{E}_t[\pi_{t+1}]=\rho_u\,\pi_t$.
  \end{itemize}
\item \emph{De dónde sale la representación}: con esas expectativas, la NKPC queda $\pi_t=\frac{\kappa}{1-\beta\rho_u}\,x_t+\frac{u_t}{1-\beta\rho_u}$, creciente en el plano (la «oferta»); y la IS con la regla de \autor{Taylor} (sin $v_t$ y con $\hat r^n_t\equiv r^n_t-\rho$),
  \begin{equation*}
  \pi_t=-\frac{(1-\rho_u)\,\theta+\phi_x}{\phi_\pi-\rho_u}\,x_t+\frac{\hat r^n_t}{\phi_\pi-\rho_u},
  \end{equation*}
  decreciente si $\phi_\pi>\rho_u$ (la «demanda»). El gráfico~\ref{fig:nek-xpi} es el caso de perturbaciones transitorias, $\rho_u=0$, en el que las expectativas son nulas.
\item \emph{Cómo se lee}: las 2 rectas relacionan las \emph{respuestas de equilibrio} de $x_t$ y $\pi_t$ a una perturbación con una persistencia dada; no son sendas por las que se mueva la economía. Si cambia la persistencia o la regla, cambian las 2 rectas: es la crítica de \autor{Lucas} aplicada al propio gráfico.
\end{itemize}
\end{anotaciones}

\begin{itemize}
\item \emph{¿Por qué un cambio de precios relativos sube el nivel general?} Con precios flexibles, que la energía se encarezca y otros bienes se abaraten no cambiaría la media. Con costes de cambiar los precios, solo los cambian las empresas con perturbaciones grandes: si unos pocos precios quieren subir mucho y muchos quieren bajar un poco, suben los primeros y los demás esperan, y el nivel de precios sube\nocite{BallMankiw1995Relative}. Por eso el encarecimiento de un insumo importante actúa como una perturbación de oferta agregada.
\item \emph{De la perturbación a la inflación, en 3 pasos}\nocite{BdE2022InformeAnual2021}:
  \begin{enumerate}
  \item \emph{efectos directos}: sube el precio de la energía que compran los hogares, que está en el IPC;
  \item \emph{efectos indirectos}: las empresas trasladan el encarecimiento de sus insumos a los precios de lo que venden;
  \item \emph{efectos de segunda ronda}: los trabajadores piden subidas de salarios para recuperar el poder adquisitivo perdido, las empresas suben los precios para defender sus márgenes y se forma la \concepto{espiral de precios y salarios}.\footnote{Trabajadores y empresas no negocian el salario real: negocian el nominal. Los trabajadores tratan de proteger su salario real pidiendo subidas nominales que compensen la inflación pasada o la que esperan, y las empresas fijan los precios sobre los costes salariales con un margen. Si los 2 quieren a la vez más de lo que hay, cada subida justifica la siguiente.}
  \end{enumerate}
  La espiral es la pescadilla que se muerde la cola: cada subida de salarios justifica la siguiente de precios, y al revés. La segunda ronda es lo que preocupa al banco central, porque convierte una perturbación de costes transitoria en inflación persistente: el tercer vértice.
\end{itemize}

% ..........................................................................
\subsubsection{Inflación inercial (\emph{built-in}): la persistencia}\label{sec:inercia}

\paragraph{a) Idea}

\begin{itemize}
\item Una vez vistos los 2 empujes, la demanda y los costes, queda el tercer vértice: lo que mantiene la inflación cuando el empuje ya ha pasado. La \concepto{persistencia} es la tendencia de la inflación a quedarse cerca de donde estaba, como un cuerpo con inercia; se mide por cuánto se parece la inflación de un trimestre a la del anterior.
\item Hay 2 fuentes de persistencia\nocite{FuhrerMoore1995Inflation}:
  \begin{itemize}
  \item la \concepto{persistencia heredada}: la inflación persiste porque persisten sus determinantes (el paro, la brecha de producción, el coste marginal);
  \item la \concepto{persistencia intrínseca}: la inflación pasada influye por sí misma en la presente (indexación, expectativas que miran atrás, contratos escalonados). Es la de las $a_i$ de~\eqref{eq:triangulo-gordon}, y recuerda a las expectativas adaptativas de los monetaristas.
  \end{itemize}
  Una tercera fuente son los retrasos de la propia política económica: el banco central tarda en reaccionar y sus decisiones tardan en hacer efecto.
\item La pregunta importa para la política: si la persistencia es intrínseca, desinflar cuesta producción aunque el banco central sea creíble.
\end{itemize}

\paragraph{b) Modelo}

\subparagraph{i) Supuestos}

\begin{itemize}
\item Partimos de la NKPC~\eqref{eq:nkpc}, con expectativas racionales y precios à la \autor{Calvo} \vertema{3.A.7}.
\item La perturbación de costes es persistente: $u_t=\rho_u\,u_{t-1}$ más una innovación de media cero, con $0\leq\rho_u<1$ ($\rho_u$, local).
\end{itemize}

\subparagraph{ii) Desarrollo}

\begin{itemize}
\item Resolviendo~\eqref{eq:nkpc} hacia delante, sustituyendo una y otra vez $\mathbb{E}_t[\pi_{t+1}]$ por su propia ecuación (con $\lim_{K\to\infty}\beta^K\mathbb{E}_t[\pi_{t+K}]=0$),
  \begin{equation}\label{eq:nkpc-adelante}
  \pi_t=\kappa\sum_{k\geq0}\beta^k\,\mathbb{E}_t[x_{t+k}]+\sum_{k\geq0}\beta^k\,\mathbb{E}_t[u_{t+k}]
  =\kappa\sum_{k\geq0}\beta^k\,\mathbb{E}_t[x_{t+k}]+\frac{u_t}{1-\beta\rho_u}.
  \end{equation}
  La inflación de hoy es el \emph{valor presente de las brechas y de los costes futuros}: las empresas que revisan hoy el precio saben que puede quedarse fijo un tiempo y miran a lo que vendrá \vertema{3.A.7}.
\item No hay ningún término de inflación pasada: toda la persistencia es heredada. Con la brecha nula, la inflación tiene la persistencia de la perturbación, $\rho_u$, y ni un ápice más. La rigidez está en el \emph{nivel} de precios, no en la tasa de inflación.
\end{itemize}

\subparagraph{iii) Implicaciones}

\begin{itemize}
\item \emph{Una desinflación creíble no costaría nada}: si el banco central anuncia de forma creíble que las brechas futuras serán las necesarias para una inflación más baja, la inflación cae de golpe, sin paro. Con precios escalonados puede incluso haber un auge durante la desinflación\nocite{Ball1994Credible}. En los datos, en cambio, la inflación baja poco a poco y con pérdidas de producción (\ref{sec:desinflacion}).
\item \emph{La curva de \autor{Phillips} neokeynesiana no es la aceleracionista}: escrita con el paro, $\pi_t=\beta\,\mathbb{E}_t[\pi_{t+1}]-a\,(U_t-U^n)+u_t$, y con el paro por debajo de su tasa natural da $\pi_t>\beta\,\mathbb{E}_t[\pi_{t+1}]$: la inflación es alta \emph{hoy} y se espera que \emph{baje}, al revés que en~\eqref{eq:aceleracionista}, de la que viene el nombre de NAIRU\nocite{Mankiw2001Inexorable}.
\end{itemize}

\paragraph{c) Extensiones}

\begin{itemize}
\item \emph{Indexación y reglas sencillas}: si los precios que no se revisan se actualizan con la inflación pasada\nocite{ChristianoEichenbaumEvans2005}, o si una parte de las empresas fija el precio con una regla sencilla que mira la inflación pasada\nocite{GaliGertler1999Inflation}, sale la \concepto{curva de \autor{Phillips} híbrida} \vertema{3.A.7}
  \begin{equation}\label{eq:nkpc-hibrida}
  \pi_t=\gamma_b\,\pi_{t-1}+\gamma_f\,\mathbb{E}_t[\pi_{t+1}]+\kappa\,x_t+u_t ,
  \end{equation}
  con $\gamma_b$ y $\gamma_f$ (locales) los pesos retrospectivo y prospectivo. Con indexación completa, $\gamma_b=1/(1+\beta)$ y $\gamma_f=\beta/(1+\beta)$, casi la mitad cada uno. Es el triángulo de \autor{Gordon} con expectativas racionales.
\item \emph{Contratos sobre salarios relativos}: si cada trabajador negocia su salario mirando a los de los contratos anteriores y a los que vendrán, la inflación hereda la inercia de los salarios\nocite{FuhrerMoore1995Inflation}.
\item \emph{Información rígida}: todos cambian el precio cada periodo, pero solo una parte de las empresas actualiza su información; las demás fijan el precio con expectativas formadas en el pasado, y las perturbaciones pasan a los precios poco a poco\nocite{MankiwReis2002Sticky}. Es la herencia de \autor{Lucas}: la información imperfecta, ahora por el coste de procesarla. Recupera la inercia sin indexación y da desinflaciones costosas.
\item \emph{Expectativas}: la última fuente de inercia son las propias expectativas de inflación. Mientras están ancladas en el objetivo, las perturbaciones pasan; si se desanclan, la inflación se perpetúa aunque la causa haya desaparecido. Por eso los bancos centrales siguen las encuestas y las expectativas que se deducen de los mercados \vertema{3.A.36}.
\item Las ecuaciones de estos modelos están en el anexo~\ref{anexo:persistencia-modelos}.
\end{itemize}

\paragraph{d) Evidencia empírica}

\begin{itemize}
\item En las estimaciones de la curva híbrida con el coste marginal real, el peso prospectivo es mayor que el retrospectivo, aunque este no es despreciable\nocite{GaliGertler1999Inflation}.
\item La persistencia no es una constante de la naturaleza: depende del régimen monetario. Es alta en los periodos sin un ancla nominal clara, como los 70, y casi desaparece con objetivos de inflación creíbles\nocite{Benati2008Investigating}. Es la crítica de \autor{Lucas} aplicada a la inflación \vertema{3.A.6}.
\end{itemize}

\paragraph{e) Valoración}

\begin{itemize}
\item La indexación arregla la persistencia, pero choca con los microdatos: implica que todos los precios cambian cada periodo, y la mitad de los precios de consumo duran más de 4 meses\nocite{BilsKlenow2004Some}. Además, con mucha indexación y pocas empresas que fijan el precio de forma óptima, estabilizar la inflación exige movimientos muy grandes de los tipos de interés.
\item La curva de \autor{Phillips} neokeynesiana sigue siendo útil para la política y para enseñar, porque con estas extensiones responde de forma gradual a las perturbaciones. Lo que queda: la inercia pesa más cuanto peor ancladas están las expectativas y más indexados están los salarios, las pensiones o los alquileres. Es la razón de la desindexación de la economía española\nocite{BOE2015Desindexacion} y de que los bancos centrales vigilen los efectos de segunda ronda.
\end{itemize}

% --------------------------------------------------------------------------
\subsection{La inflación de 2021-2026 a la luz del triángulo}\label{sec:inflacion-reciente}

\begin{itemize}
\item Una vez vistos los 3 vértices, los aplicamos al episodio que ha devuelto la inflación al centro del debate. En España el IPC llegó al 10,8~\% en julio de 2022 (el IAPC, al 10,7~\%)\nocite{INE2026IPCSerie,Eurostat2026IAPC}, y la inflación de la zona del euro, al 10,6~\% en octubre de ese año\nocite{Eurostat2026IAPC}. La descomposición sigue los 3 bloques \vertema{4.A.29}.
\end{itemize}

\subsubsection{Costes}\label{sec:reciente-costes}

\begin{itemize}
\item La reapertura tras la pandemia chocó con cuellos de botella en las cadenas de suministro, y la invasión de Ucrania disparó el gas, la electricidad y los alimentos \vertema{4.A.5}. En la zona del euro, la energía llegó a aportar 3,8 puntos a la inflación en 2022, los alimentos 2,9 y los cuellos de botella 1,4 a comienzos de 2023, mientras que el mercado de trabajo apenas aportó 0,2\nocite{ArceEtAl2024What,ArceEtAl2024SUERF}.
\item Las empresas trasladaron los costes y, en algunos sectores, ampliaron márgenes: entre el primer trimestre de 2022 y el primero de 2023, los beneficios unitarios explican un 45~\% del aumento del deflactor del consumo y los precios de importación, un 40~\%\nocite{HansenToscaniZhou2023}.
\end{itemize}

\subsubsection{Demanda}\label{sec:reciente-demanda}

\begin{itemize}
\item En Estados Unidos, el estímulo fiscal de 2021, mucho mayor que la brecha de producción, y el ahorro acumulado en la pandemia empujaron la demanda\nocite{Summers2021Biden,Blanchard2021Defense}, y la curva de \autor{Phillips}, que en los 2010 parecía plana, se empinó con un mercado de trabajo muy tenso\nocite{Blanchard2016Phillips,BenignoEggertsson2023}. Con la inflación alta, las empresas cambiaron los precios más a menudo\nocite{GautierEtAl2026WP3181}.
\item En la zona del euro, la demanda pesó menos al principio: la subida de la subyacente de la segunda mitad de 2021 fue sobre todo de oferta, y la demanda fue ganando peso hasta igualarla a mediados de 2022\nocite{GoncalvesKoester2022}.
\end{itemize}

\subsubsection{Inercia}\label{sec:reciente-inercia}

\begin{itemize}
\item \textcite{BernankeBlanchard2025What} llegan para Estados Unidos a la misma conclusión que \textcite{ArceEtAl2024What} para la zona del euro: la inflación nació de perturbaciones de precios, y la tensión del mercado de trabajo, que llegó más tarde, es lo que la hizo persistir. Las expectativas a largo plazo se mantuvieron cerca del 2~\%, y la inflación bajó sin una recesión.\footnote{\autor{Ben Bernanke} compartió el Premio Nobel de 2022 con \autor{Douglas Diamond} y \autor{Philip Dybvig} «por sus investigaciones sobre los bancos y las crisis financieras»\nocite{Nobel2022}.}
\item \emph{2026}: la energía vuelve a encarecerse y la inflación general sube, pero la subyacente no (gráfico~\ref{fig:inflacion-general-subyacente}): en octubre, la de la zona del euro está en el 2,5~\%\nocite{Eurostat2026IAPCoct} y la de España, la del INE, en el 3,1~\%\nocite{INE2026IPCAdelantadoOct}; en Estados Unidos, el IPC sube un 3,7~\% en septiembre\nocite{BLS2026IPC}. Con la lección de 2022, el BCE subió los tipos en junio y en septiembre (la facilidad de depósito, al 2,25~\% y al 2,50~\%) para que no se desanclen las expectativas\nocite{BCE2026Tipos}, que siguen en el 2,0~\%\nocite{BCE2026SPF4T} \vertema{3.A.36}.
\end{itemize}

\begin{figure}[H]
\centering
\caption{Inflación general y subyacente en la zona del euro y en España, 1999-2026}\label{fig:inflacion-general-subyacente}
\grafico[0.85\textwidth]{inflacion-general-subyacente}
\fuente{Elaboración propia con datos de Eurostat (IAPC, prc\_hicp\_minr) y, para la subyacente anterior a 2001-2002, del BCE\nocite{Eurostat2026IAPC,BCE2026ICP}. Tasa interanual, en~\%; subyacente: sin energía ni alimentos (la de España con la definición de Eurostat, no la del INE). El dato de octubre de 2026 es un avance. Qué hay que ver: en 2022 la perturbación de costes pasó a la subyacente con retraso (la segunda ronda); en 2026, de momento, no.}
\end{figure}

\begin{anotaciones}
\textbf{¿Oferta o demanda? Una descomposición con datos} (para consultar; en el cante, una cifra):
\begin{itemize}
\item \emph{El método}\nocite{Shapiro2022Oferta}: para cada uno de los componentes del deflactor del consumo, se estiman con un modelo las sorpresas de su precio y de su cantidad. Si las 2 van en el mismo sentido (sube el precio y sube la cantidad), el componente se clasifica ese mes como de \emph{demanda}; si van en sentido contrario (sube el precio y baja la cantidad), como de \emph{oferta}; si no son significativas, como \emph{ambiguo}. Las contribuciones se suman con los pesos de cada componente.
\item \emph{Los resultados en Estados Unidos} (gráfico~\ref{fig:inflacion-oferta-demanda}): en junio de 2022, de un 7,1~\% de inflación interanual del deflactor del consumo, 3,4 puntos eran de oferta, 2,8 de demanda y 1,0 ambiguos; en septiembre de 2026, de un 3,7~\%, 1,3 puntos son de oferta, 1,2 de demanda y 1,2 ambiguos\nocite{FRBSF2026OfertaDemanda}.
\item \emph{En la zona del euro}, el BCE aplicó el método a la inflación sin energía ni alimentos (usando la cifra de negocios como cantidad): la subida de la segunda mitad de 2021 fue sobre todo de oferta, y a mediados de 2022 oferta y demanda aportaban ya parecido; los bienes industriales, de oferta, y los servicios, de demanda\nocite{GoncalvesKoester2022}.
\item \emph{Límite}: la clasificación es binaria (cada componente es todo de oferta o todo de demanda en cada mes) y no dice cuánto pesa cada causa dentro de un componente.
\end{itemize}
\begin{figuracaja}
\caption{Inflación de oferta y de demanda en Estados Unidos, 2019-2026}\label{fig:inflacion-oferta-demanda}
\grafico[0.95\textwidth]{inflacion-oferta-demanda}
\fuente{Elaboración propia con datos del Banco de la Reserva Federal de San Francisco\nocite{FRBSF2026OfertaDemanda}, con el método de \textcite{Shapiro2022Oferta}. Deflactor del consumo privado (PCE), tasa interanual y contribuciones, en puntos porcentuales; enero de 2019 a septiembre de 2026.}
\end{figuracaja}
\end{anotaciones}

\begin{itemize}
\item \emph{Recapitulación}: el dinero explica la inflación a largo plazo; el triángulo de \autor{Gordon}, la de cada año, con sus 3 vértices; y la inercia, por qué cuesta bajarla. La inflación tiene 3 orígenes, pero solo se queda si el dinero, el déficit o las expectativas la alimentan. Una vez vistas sus causas, pasamos a lo que cuesta.
\end{itemize}

% ==========================================================================
\section{Efectos sobre la eficiencia económica y el bienestar}\label{sec:efectos}
\minutos{8}

\begin{itemize}
\item Una vez vistas las causas, pasamos a los efectos, la otra mitad del título: qué cuesta la inflación, a quién y cuánta conviene.
\item \emph{Conclusión anticipada}: la inflación cuesta aunque se anticipe, porque encarece el dinero, obliga a cambiar precios y distorsiona los impuestos; si no se anticipa, además redistribuye y crea incertidumbre. La tasa óptima depende de las fricciones: negativa con solo la del dinero (la regla de \autor{Friedman}), cero con precios rígidos y algo positiva con el límite inferior de los tipos y los salarios rígidos a la baja. De ahí el 2~\%.
\end{itemize}

% --------------------------------------------------------------------------
\subsection{Costes de la inflación anticipada y no anticipada}\label{sec:costes-inflacion}

\begin{itemize}
\item La distinción clave es si la inflación se \emph{anticipa}. Los contratos fijan el tipo nominal con la inflación esperada, $i\simeq r+\pi^e$ (ecuación de \autor{Fisher}) \vertema{3.A.3}, y el tipo real que de verdad se obtiene es $i-\pi$:
  \begin{itemize}
  \item con \concepto{inflación anticipada} ($\pi=\pi^e$), todos la incorporan a sus contratos y nadie gana ni pierde por sorpresa; cuesta por lo que la economía hace para protegerse de ella;
  \item con \concepto{inflación no anticipada} ($\pi\neq\pi^e$), el tipo real efectivo se aparta del pactado y la riqueza pasa de unos a otros.
  \end{itemize}
\item Los costes de la anticipada son sobre todo de \emph{eficiencia}; los de la no anticipada, de \emph{bienestar} y de distribución. La frontera es borrosa: la inflación alta es casi siempre variable y, por tanto, mal anticipada\nocite{Friedman1977Nobel}.
\end{itemize}

\subsubsection{Costes de la inflación anticipada}\label{sec:anticipada}

\paragraph{Los costes en saldos reales y el triángulo de \autor{Bailey}}

\begin{itemize}
\item El dinero no rinde intereses: tenerlo cuesta el tipo nominal, $i=r+\pi^e$. Con más inflación el público tiene menos saldos reales y, para pagar lo mismo, va más a menudo al banco a sacar efectivo y dedica tiempo y recursos a gestionar su liquidez \vertema{3.A.35}. Se llaman \concepto{costes en suela de zapatos} (\emph{shoe-leather costs}) porque, en el modelo de inventarios de \autor{Baumol} y \autor{Tobin}, el que tiene menos efectivo tiene que ir más veces al banco, y las suelas se gastan en los viajes: la imagen es literal.\footnote{\autor{James Tobin} recibió el Premio Nobel en 1981 «por su análisis de los mercados financieros y de sus relaciones con las decisiones de gasto, el empleo, la producción y los precios»\nocite{Nobel1981Tobin}.} Con inflaciones altas, la imagen se queda corta: las empresas pagan los salarios varias veces por semana y la gente corre a gastar el dinero en cuanto lo cobra.
\item \textcite{Bailey1956Welfare} los mide como una pérdida de excedente. Sea $L(i)\equiv M^D/P$ la demanda de saldos reales, decreciente en el tipo nominal. Como cualquier curva de demanda, su altura en cada punto es lo que el público está dispuesto a pagar por la última unidad de saldos reales, es decir, el valor de los servicios que le presta (ahorrar viajes al banco, comprar sin fricciones). Producir dinero no cuesta casi nada, así que lo eficiente es saciar la demanda, $L(0)$. Con $i>0$ el público renuncia a los saldos entre $L(i)$ y $L(0)$, que valoraba en el área bajo la demanda; de ese área, el rectángulo $i\,L(i)$ no se pierde: es el impuesto inflacionario, una transferencia al emisor. Lo que queda es el triángulo de la pérdida (gráfico~\ref{fig:bailey}).
\end{itemize}

\begin{figure}[H]
\centering
\caption{El coste de bienestar de la inflación: el triángulo de \autor{Bailey}}\label{fig:bailey}
\grafico[0.6\textwidth]{bailey}
\fuente{Elaboración propia a partir de \textcite{Bailey1956Welfare} y \textcite{Lucas2000Inflation}. Con el tipo nominal $i_0=r+\pi$ el público tiene $L(i_0)$; el rectángulo es el impuesto inflacionario (una transferencia al emisor) y el triángulo sombreado, la pérdida de eficiencia, que desaparece con $i=0$.}
\end{figure}

\begin{itemize}
\item \emph{De dónde sale la función}: el área bajo la demanda entre $0$ y $i$, medida sobre el eje del tipo, es $\int_0^{i}L(x)\,dx$ (el excedente que se pierde al pasar de $i=0$ a $i$: el rectángulo del impuesto más el triángulo); restándole el rectángulo, que es transferencia, la \concepto{pérdida de \autor{Bailey}} es
  \funcionclave[eq:bailey]{\mathrm{PB}(i)=\int_0^{i}L(x)\,dx-i\,L(i)\;\approx\;\tfrac12\,\lvert L'(0)\rvert\,i^2}
  con $\mathrm{PB}(i)$ la pérdida de bienestar por periodo (en redonda; $L$, de liquidez). La aproximación es un desarrollo de segundo orden alrededor de $i=0$: $\mathrm{PB}(0)=0$, $\mathrm{PB}'(i)=L(i)-L(i)-i\,L'(i)=-i\,L'(i)$, que es nula en $i=0$, y $\mathrm{PB}''(0)=-L'(0)>0$. Geométricamente, es el área de un triángulo de altura $i$ y base $\lvert L'(0)\rvert\,i$.
\item \emph{Lectura}: la pérdida crece con el \emph{cuadrado} del tipo nominal, como el exceso de gravamen de cualquier impuesto: pasar del 2 al 4~\% cuesta 3 veces más que pasar del 0 al 2~\%. La inflación baja cuesta poco; la alta, mucho. Con la demanda de \autor{Cagan}, $L(i)=\zeta e^{-\alpha i}$, $\lvert L'(0)\rvert=\alpha\zeta$: la misma semielasticidad $\alpha$ que fija el señoreaje máximo (\ref{sec:cagan}).
\item \emph{Evidencia}: con datos de Estados Unidos de 1900 a 1994, pasar de una inflación del 10~\% a una del 0~\% equivale a un aumento de la renta real de algo menos del 1~\%\nocite{Lucas2000Inflation}. Pequeño con inflación baja, grande con inflación alta.
\end{itemize}

\paragraph{Costes de menú y dispersión de los precios relativos}

\begin{itemize}
\item Cambiar los precios cuesta (rehacer catálogos, avisar a los clientes, decidir): son los \concepto{costes de menú}. Pequeños para cada empresa, bastan para que los precios no se ajusten a la vez\nocite{Mankiw1985Small} \vertema{3.A.7}.
\item El coste mayor es indirecto: como los precios no se revisan a la vez, con inflación los precios relativos de bienes iguales se separan, y los hogares compran cantidades distintas de variedades que valoran igual. Por eso en la nueva economía keynesiana la inflación entra en la pérdida social con un peso grande, y la estabilidad de precios es el óptimo con precios rígidos\nocite{Woodford2003Interest} \vertema{3.A.7}.
\end{itemize}

\paragraph{La unidad de cuenta y la información de los precios}

\begin{itemize}
\item Con inflación el dinero mide peor: los contratos a largo plazo, la contabilidad de las empresas y las comparaciones de precios en el tiempo se complican.
\item Si además la inflación es variable, una subida de un precio deja de distinguir un cambio de los precios relativos de uno general: la señal de los precios se debilita y los recursos se asignan peor\nocite{Lucas1973Some,Friedman1977Nobel}.
\end{itemize}

\paragraph{Las distorsiones fiscales}

\begin{itemize}
\item Los impuestos se calculan sobre magnitudes nominales, y la inflación, aunque se anticipe, cambia lo que de verdad se paga:
  \begin{itemize}
  \item \emph{progresividad en frío}: si los tramos del impuesto sobre la renta no se actualizan, una subida del salario nominal que solo compensa los precios lleva a tramos más altos, y el tipo medio sube sin que suba la renta real;
  \item \emph{rentas del capital}: se grava el interés nominal. Con un tipo impositivo $\tau$, el rendimiento real después de impuestos es $(1-\tau)\,i-\pi=(1-\tau)\,r-\tau\,\pi$: la inflación lo reduce en $\tau\pi$. Con $r=2$~\%, $\tau=30$~\% y $\pi=2$~\%, baja del 1,4~\% al 0,8~\%. Lo mismo hacen las amortizaciones a coste histórico en el impuesto de sociedades;
  \item por esta interacción con los impuestos sobre el ahorro y la vivienda, una inflación del 2~\% cuesta en Estados Unidos alrededor del 1~\% del PIB \emph{cada año}, frente a un coste único de desinflar de alrededor del 5~\%\nocite{Feldstein1997Costs}: mucho más que la suela de zapatos.
  \end{itemize}
\item \emph{Efecto \autor{Olivera}-\autor{Tanzi}}: entre el hecho imponible y el cobro pasa tiempo, y la recaudación pierde valor real. Con un retraso de 2 meses y una inflación del 10~\% mensual se pierde un 17~\% ($1/1{,}1^2\simeq0{,}83$). Con inflación alta el déficit crece, hay que emitir más dinero y la inflación se alimenta a sí misma\nocite{Olivera1967Money,Tanzi1977Inflation}: es un eslabón de las hiperinflaciones (\ref{sec:hiperinflacion}).
\end{itemize}

\paragraph{Inflación y acumulación de capital}

\begin{itemize}
\item ¿Cambia la inflación anticipada el capital a largo plazo? Las respuestas dependen de por qué se tiene dinero:
  \begin{itemize}
  \item el \concepto{efecto \autor{Mundell}-\autor{Tobin}}: si el dinero es un activo alternativo al capital, más inflación desplaza la cartera hacia el capital, baja el tipo real y sube el capital de estado estacionario\nocite{Mundell1963Inflation,Tobin1965} \vertema{3.A.43};\footnote{\autor{Robert Mundell} recibió el Premio Nobel en 1999 «por su análisis de la política monetaria y fiscal con distintos regímenes de tipo de cambio y por su análisis de las áreas monetarias óptimas»\nocite{Nobel1999Mundell}.}
  \item la \concepto{superneutralidad}: con hogares que optimizan a lo largo de su vida y el dinero en la utilidad (MIU), el capital de estado estacionario lo fija la regla de oro modificada y no depende de la inflación\nocite{Sidrauski1967Rational} \vertema{3.A.35};
  \item si el dinero hace falta también para invertir (restricción de pago por adelantado), la inflación grava la inversión y el capital \emph{baja}\nocite{Stockman1981Anticipated}.
  \end{itemize}
\item Son resultados de estado estacionario; en los datos, la inflación alta va con menos inversión y menos crecimiento, sobre todo por su volatilidad.
\end{itemize}

\subsubsection{Costes de la inflación no anticipada}\label{sec:no-anticipada}

\begin{itemize}
\item \emph{Redistribución}: una inflación mayor que la esperada baja el tipo real \emph{ex post} y pasa riqueza de los acreedores a los deudores. En Estados Unidos pierden sobre todo los hogares mayores y ricos, que tienen bonos, y ganan los jóvenes de clase media con hipotecas a tipo fijo y el sector público\nocite{DoepkeSchneider2006Inflation}. Es una redistribución que nadie ha votado. Con deuda pública alta, la inflación inesperada y los tipos reales negativos la reducen, como tras la Segunda Guerra Mundial\nocite{ReinhartSbrancia2015Liquidation} \vertema{3.A.39}.
\item \emph{Regresividad}: los hogares de renta baja tienen una parte mayor de su riqueza en efectivo y gastan más en energía y alimentos, así que el impuesto inflacionario y las subidas de 2022 les pesaron más\nocite{BCE2022RentaBaja}.
\item \emph{Incertidumbre}: la inflación alta es más variable; los contratos se acortan, los tipos a largo plazo llevan una prima por el riesgo de inflación y la inversión se retrae. \textcite{Friedman1977Nobel} sostuvo que por esto, a medio plazo, más inflación puede ir con \emph{más} paro, no con menos.
\item \emph{Competitividad}: con tipo de cambio fijo o en una unión monetaria, una inflación mayor que la de los socios aprecia el tipo de cambio real y resta competitividad \vertema{3.B.11}.
\end{itemize}

\subsubsection{¿Tiene la inflación algún beneficio?}\label{sec:beneficios}

\begin{itemize}
\item Hay consenso en un \emph{objetivo} de inflación positivo y bajo, no en que la inflación sea buena para la eficiencia. Lo que tiene de cierto:
  \begin{itemize}
  \item \emph{lubricante del mercado de trabajo}: si los salarios nominales no bajan, una inflación moderada permite que bajen los reales donde hace falta\nocite{Tobin1972Inflation,AkerlofDickensPerry1996} \vertema{3.A.28};
  \item \emph{ingreso público}: el señoreaje, para un Estado con un sistema fiscal pobre\nocite{Phelps1973Inflation} \vertema{3.A.39};
  \item \emph{margen para la política monetaria}: con una inflación media positiva, el tipo nominal medio es mayor y queda más espacio para bajarlo antes de tocar el límite inferior \vertema{3.A.37}.
  \end{itemize}
\item Ninguno justifica una inflación alta, pero los 3 explican por qué el objetivo no es cero: es la pregunta de la tasa óptima. La tabla~\ref{tab:costes-inflacion} resume los costes y los beneficios.
\end{itemize}

\begin{table}[H]
\centering
\caption{Los costes y los beneficios de la inflación}\label{tab:costes-inflacion}
\small
\begin{tabular}{>{\raggedright\arraybackslash}m{3.1cm}>{\centering\arraybackslash}m{2.3cm}>{\raggedright\arraybackslash}m{5.2cm}>{\raggedright\arraybackslash}m{3.4cm}}
\toprule
\textbf{Efecto} & \textbf{Inflación} & \textbf{Mecanismo} & \textbf{Referencia}\\
\midrule
\rowcolor{tceeazulpalido}\multicolumn{4}{l}{\emph{Eficiencia}}\\
Suela de zapatos & anticipada & menos saldos reales; triángulo de pérdida & \autor{Bailey} (1956), \autor{Lucas} (2000)\\
Menú y dispersión & anticipada & precios relativos distorsionados & \autor{Mankiw} (1985), \autor{Woodford} (2003)\\
Fiscales & anticipada & impuestos sobre lo nominal; retraso del cobro & \autor{Feldstein} (1997), \autor{Tanzi} (1977)\\
Información & sobre todo no anticipada & los precios señalan peor & \autor{Lucas} (1973), \autor{Friedman} (1977)\\
\rowcolor{tceesalmon}\multicolumn{4}{l}{\emph{Bienestar y distribución}}\\
Redistribución & no anticipada & de acreedores a deudores & \autor{Doepke} y \autor{Schneider} (2006)\\
Regresividad & las 2 & más efectivo y más energía en la cesta & BCE (2022)\\
Incertidumbre & no anticipada & contratos cortos, prima de riesgo & \autor{Friedman} (1977)\\
\rowcolor{tceesalvia}\multicolumn{4}{l}{\emph{Beneficios}}\\
Lubricante del mercado de trabajo & anticipada & baja el salario real donde el nominal no baja & \autor{Tobin} (1972); \autor{Akerlof}, \autor{Dickens} y \autor{Perry} (1996)\\
Ingreso público & anticipada & señoreaje con un sistema fiscal pobre & \autor{Phelps} (1973)\\
Margen para la política monetaria & anticipada & tipo nominal medio más alto, lejos del límite inferior & \autor{Blanchard}, \autor{Dell'Ariccia} y \autor{Mauro} (2010)\\
\bottomrule
\end{tabular}
\fuente{Elaboración propia a partir de las obras citadas.}
\end{table}

% --------------------------------------------------------------------------
\subsection{¿Existe una tasa de inflación óptima?}\label{sec:inflacion-optima}

\begin{itemize}
\item Vistos los costes, la pregunta normativa es cuánta inflación conviene. La teoría da 3 respuestas: deflación, estabilidad de precios o algo de inflación\nocite{SchmittGroheUribe2010Optimal}.
\end{itemize}

\subsubsection{La regla de \autor{Friedman}}\label{sec:regla-friedman}

\paragraph{a) Idea}

\begin{itemize}
\item \textcite{Friedman1969Optimum} parte del triángulo de \autor{Bailey}: el coste privado de tener dinero es el tipo nominal y su coste social de producción es casi nulo. La eficiencia exige igualarlos.
\end{itemize}

\paragraph{b) Modelo}

\subparagraph{i) Supuestos}

\begin{itemize}
\item Precios flexibles, un dinero que presta un servicio (da utilidad o hace falta para comprar) y cuya producción no cuesta, y un Estado que puede recaudar con impuestos de suma fija.
\end{itemize}

\subparagraph{ii) Desarrollo}

\begin{itemize}
\item En el modelo con dinero en la utilidad (MIU, \emph{money in the utility function})\nocite{Sidrauski1967Rational}, el hogar demanda saldos reales hasta que su relación marginal de sustitución con el consumo iguala su coste de oportunidad, $u_m/u_c=i$ (con tiempo discreto, $i/(1+i)$; la derivación, \vertema{3.A.35}). Lo eficiente es saciar la demanda, $u_m=0$, es decir,
  \begin{equation}\label{eq:regla-friedman}
  i=r+\pi=0\quad\Longrightarrow\quad\pi^*=-r .
  \end{equation}
  Es la \concepto{regla de \autor{Friedman}}: una deflación igual al tipo de interés real. Con ella la pérdida~\eqref{eq:bailey} desaparece (gráfico~\ref{fig:bailey}).\footnote{No es la regla del $k$~\%, un crecimiento constante del dinero, $\gamma_M=\overline\gamma_M$, que el mismo \autor{Friedman} propuso en 1960 para estabilizar \vertema{3.A.5}, ni el dinero lanzado desde un helicóptero, que es una imagen del mismo libro de 1969 \vertema{3.A.36}.}
\end{itemize}

\subparagraph{iii) Implicaciones}

\begin{itemize}
\item Con un tipo real del 2 al 4~\%, la inflación óptima estaría entre el $-2$ y el $-4$~\%\nocite{SchmittGroheUribe2010Optimal}. Ningún banco central la sigue: los objetivos de los países avanzados están en torno al 2~\%.
\end{itemize}

\paragraph{c) Extensiones: la inflación como impuesto}

\begin{itemize}
\item \emph{Sin impuestos de suma fija}, la inflación es un impuesto más, y la imposición óptima pide usarlo un poco, $\pi^*>-r$\nocite{Phelps1973Inflation} \vertema{3.A.39}; pero si el dinero es un bien intermedio, la regla de \autor{Friedman} vuelve a ser óptima\nocite{ChariChristianoKehoe1996Optimality}. Las rentas que escapan al fisco y la demanda exterior de la moneda empujan en el mismo sentido, con poca fuerza. El desarrollo está en el anexo~\ref{anexo:inflacion-optima}.
\end{itemize}

\paragraph{d) Evidencia empírica}

\begin{itemize}
\item La ganancia de ir de una inflación baja a la regla es pequeña y difícil de medir\nocite{Lucas2000Inflation}. La única deflación duradera de un país avanzado en la posguerra, la de Japón, no fue una elección y vino con estancamiento (\ref{sec:deflacion}).
\end{itemize}

\paragraph{e) Valoración: las rigideces nominales}

\begin{itemize}
\item La regla ignora que los precios son rígidos. Si no se revisan a la vez, cualquier inflación, también la negativa, dispersa los precios relativos, y la inflación óptima es \emph{cero}: la \concepto{estabilidad de precios}\nocite{Woodford2003Interest}. Con las 2 fricciones, la del dinero y la de los precios, el óptimo queda entre $-r$ y 0, mucho más cerca de la estabilidad de precios\nocite{SchmittGroheUribe2010Optimal}. El desarrollo, en el anexo~\ref{anexo:inflacion-optima}.
\item La teoría da, pues, una inflación óptima entre $-r$ y algo apenas por debajo de cero. La práctica ha elegido el 2~\%: hay que explicar la diferencia.
\end{itemize}

\subsubsection{¿Por qué el 2~\%?}\label{sec:dos-por-ciento}

\begin{itemize}
\item Hay 3 razones para un objetivo por encima de lo que dice la teoría:
  \begin{itemize}
  \item \emph{el sesgo de medición}: el IPC sobrestima la inflación verdadera porque recoge mal las mejoras de calidad, los productos nuevos y la sustitución entre bienes; en Estados Unidos, a mediados de los 90, 1,1 puntos al año, de los que 0,6 eran de calidad\nocite{Boskin1996Toward} \vertema{3.A.9}. Pero corregir el objetivo solo es óptimo si los precios rígidos son los ajustados por calidad\nocite{SchmittGroheUribe2010Optimal};
  \item \emph{el límite inferior de los tipos}: con una inflación media más alta, el tipo nominal medio es más alto y queda más margen para bajarlo en una recesión \vertema{3.A.37}. En un modelo estimado para Estados Unidos antes de 2008, el riesgo era pequeño (el tipo nominal óptimo tenía una media del 4,4~\% y una desviación típica del 0,9~\%, y la inflación óptima, una media del $-0{,}4$~\%)\nocite{SchmittGroheUribe2010Optimal}; la caída del tipo natural y los años en el límite tras 2008 lo agrandaron, y se llegó a proponer estudiar un objetivo del 4~\%\nocite{BlanchardDellAricciaMauro2010} \vertema{3.A.36} \vertema{3.A.40};
  \item \emph{la rigidez de los salarios a la baja}: si los salarios nominales no bajan, un poco de inflación facilita los ajustes de los salarios reales\nocite{Tobin1972Inflation,AkerlofDickensPerry1996}; los modelos estimados piden por esto menos de medio punto\nocite{SchmittGroheUribe2010Optimal}.
  \end{itemize}
  Fijar el 2~\% es curarse en salud: un seguro contra el límite inferior, contra los errores de medición y contra la rigidez de los salarios.
\item En una unión monetaria hay una razón más: si la inflación media es baja, los países que necesitan abaratarse en términos reales tendrían que pasar por la deflación.
\item La tabla~\ref{tab:inflacion-optima} resume la tasa óptima según las fricciones.
\item \emph{En la práctica}: el BCE persigue desde julio de 2021 un objetivo \emph{simétrico} del 2~\% a medio plazo, confirmado en junio de 2025: las desviaciones a la baja preocupan tanto como al alza\nocite{BCE2021Estrategia,BCE2025Estrategia} \vertema{3.A.36}.\footnote{El Banco de España explica la estrategia y los instrumentos con que el BCE responde a las desviaciones en las 2 direcciones\nocite{BdEEstrategiaBCE,BdEInstrumentosBCE}.} La Reserva Federal tiene un objetivo del 2~\% con el deflactor del consumo.\footnote{En agosto de 2020 la Reserva Federal adoptó un objetivo de inflación media, que compensaba los años por debajo del 2~\% con otros por encima\nocite{Fed2020Prensa}; tras la inflación de 2021-2023, en agosto de 2025 lo abandonó y volvió a un objetivo flexible de inflación, sin compensación\nocite{Fed2025Prensa,Fed2025Revision,Powell2025Jackson}. Compensar, como un objetivo de nivel de precios, ayuda en el límite inferior, porque el público espera más inflación cuando el tipo nominal no puede bajar\nocite{EggertssonWoodford2003Zero}, pero exige mucha credibilidad \vertema{3.A.36}.}
\end{itemize}

\begin{table}[htbp]
\centering
\caption{La tasa de inflación óptima según las fricciones}\label{tab:inflacion-optima}
\small
\begin{tabular}{>{\raggedright\arraybackslash}m{4.4cm}>{\centering\arraybackslash}m{3.6cm}>{\raggedright\arraybackslash}m{5.4cm}}
\toprule
\textbf{Fricción que se añade} & \textbf{Inflación óptima} & \textbf{Autores}\\
\midrule
\rowcolor{tceeazulpalido}\multicolumn{3}{l}{\emph{Teoría}}\\
Solo la monetaria (tener dinero cuesta $i$) & $\pi^*=-r$ & \autor{Friedman} (1969)\\
Impuestos distorsionantes & algo mayor que $-r$; $-r$ si el dinero es un bien intermedio & \autor{Phelps} (1973); \autor{Chari}, \autor{Christiano} y \autor{Kehoe} (1996)\\
Precios rígidos & $\approx0$ & \autor{Woodford} (2003); \autor{Schmitt-Grohé} y \autor{Uribe} (2010)\\
\rowcolor{tceesalmon}\multicolumn{3}{l}{\emph{Práctica}}\\
Límite inferior, salarios rígidos a la baja, sesgo de medición & algo positiva & \autor{Akerlof}, \autor{Dickens} y \autor{Perry} (1996); \autor{Boskin} \emph{et al.} (1996); \autor{Blanchard}, \autor{Dell'Ariccia} y \autor{Mauro} (2010)\\
Objetivo del BCE & 2~\% simétrico a medio plazo & BCE (2021, 2025)\\
\bottomrule
\end{tabular}
\fuente{Elaboración propia a partir de \textcite{SchmittGroheUribe2010Optimal} y de las obras citadas.}
\end{table}

% --------------------------------------------------------------------------
\subsection{Inflación y estabilización macroeconómica}\label{sec:estabilizacion}

\begin{itemize}
\item Además del nivel medio importa cómo se mueve la inflación con el ciclo. Con perturbaciones de demanda, estabilizar la inflación estabiliza también la actividad (la divina coincidencia); con perturbaciones de costes hay que elegir, y el banco central reparte el ajuste entre inflación y producción según el peso que dé a cada una (gráfico~\ref{fig:nek-xpi}) \vertema{3.A.7} \vertema{3.A.36}.
\item Un objetivo de inflación creíble es lo que permite esa flexibilidad: con las expectativas ancladas, una perturbación de costes pasa menos a la inflación esperada, y el banco central puede mirar a través de ella sin que se convierta en inercia. Es lo que se vio en 2022-2024 y lo que se pone a prueba en 2026 (\ref{sec:inflacion-reciente}).
\item \emph{Recapitulación}: la inflación cuesta en eficiencia aunque se anticipe y en bienestar cuando no; la teoría pide entre $-r$ y 0, y la práctica, el 2~\%. En los extremos esos costes se disparan: lo vemos ahora.
\end{itemize}

% ==========================================================================
\section{Contextos extremos: hiperinflación y deflación}\label{sec:extremos}
\minutos{5,5}

\begin{itemize}
\item Una vez vistos los costes de la inflación moderada, pasamos a los 2 extremos, donde se multiplican.
\item \emph{Conclusión anticipada}: la hiperinflación es fiscal en su origen y monetaria en su mecanismo, y se acaba de golpe con una reforma fiscal creíble; la deflación es buena si viene de la oferta y peligrosa si viene de una demanda hundida con deudas altas, porque encarece las deudas y el tipo real justo cuando la política monetaria se queda sin margen.
\end{itemize}

% --------------------------------------------------------------------------
\subsection{Hiperinflación}\label{sec:hiperinflacion}

\subsubsection{Idea}\label{sec:hiperinflacion-idea}

\begin{itemize}
\item Siguiendo a \textcite{Cagan1956Monetary}, hay \concepto{hiperinflación} cuando los precios suben más de un 50~\% al mes, casi un 13.000~\% al año ($1{,}5^{12}-1\simeq128{,}7$). Ha habido 56 hasta 2012, casi todas en el siglo \textsc{xx}, tras guerras, revoluciones o el hundimiento de un régimen, y Venezuela es la 57.ª (tabla~\ref{tab:hiperinflaciones})\nocite{HankeKrus2013World,HankeBushnell2017Measuring}.
\item Detrás de todas hay 2 cosas: un Estado que no puede financiarse con impuestos ni con deuda y recurre al banco central, y una pérdida de confianza en el dinero que hace que el público se deshaga de él.\footnote{En 1921 el calendario de Londres exigió a Alemania 2.000 millones de marcos oro al año más el 26~\% del valor de sus exportaciones en concepto de reparaciones\nocite{HansardReparaciones1921}. El Gobierno financió su déficit con emisión, y en octubre de 1923 los precios subían un 29.500~\% al mes\nocite{HankeKrus2013World}; la hiperinflación terminó con la reforma monetaria y fiscal de finales de 1923\nocite{Sargent1982Ends}. La crónica completa, en \href{https://es.wikipedia.org/wiki/Hiperinflaci\%C3\%B3n_en_la_Rep\%C3\%BAblica_de_Weimar}{«Hiperinflación en la República de Weimar», Wikipedia}.}
\end{itemize}

\begin{table}[H]
\centering
\caption{Algunas hiperinflaciones}\label{tab:hiperinflaciones}
\small
\begin{tabular}{>{\raggedright\arraybackslash}m{2.1cm}>{\centering\arraybackslash}m{3.4cm}>{\centering\arraybackslash}m{2.2cm}>{\centering\arraybackslash}m{3.1cm}>{\centering\arraybackslash}m{2.6cm}}
\toprule
\textbf{País} & \textbf{Periodo} & \textbf{Mes del máximo} & \textbf{Inflación mensual máxima} & \textbf{Los precios se duplicaban en}\\
\midrule
Hungría & ago. 1945-jul. 1946 & jul. 1946 & $4{,}19\cdot10^{16}$~\% & 15 horas\\
Zimbabue & mar. 2007-nov. 2008 & nov. 2008 & $7{,}96\cdot10^{10}$~\% & 24,7 horas\\
Yugoslavia & abr. 1992-ene. 1994 & ene. 1994 & $3{,}13\cdot10^{8}$~\% & 1,4 días\\
Alemania & ago. 1922-dic. 1923 & oct. 1923 & 29.500~\% & 3,7 días\\
Grecia & may. 1941-dic. 1945 & oct. 1944 & 13.800~\% & 4,3 días\\
China & oct. 1947-may. 1949 & abr. 1949 & 5.070~\% & 5,3 días\\
Venezuela & desde nov. 2016 & — & — & —\\
\bottomrule
\end{tabular}
\fuente{\textcite{HankeKrus2013World}; Venezuela, \textcite{HankeBushnell2017Measuring} y FMI\nocite{FMI2026InflacionFinAno}. Para Venezuela no hay una tasa mensual máxima comparable; la inflación anual a fin de año fue del 862,6~\% en 2017, el 130.060~\% en 2018 y el 9.586~\% en 2019.}
\end{table}

\subsubsection{Señoreaje e impuesto inflacionario}\label{sec:senoreaje-extremos}

\begin{itemize}
\item Emitir dinero da al Estado un ingreso real, el \concepto{señoreaje} \vertema{3.A.39}:
  \begin{equation}\label{eq:senoreaje}
  S_t\equiv\frac{M_t-M_{t-1}}{P_t}=\underbrace{\left(\frac{M_t}{P_t}-\frac{M_{t-1}}{P_{t-1}}\right)}_{\text{aumento de los saldos reales}}+\underbrace{\frac{\pi_t}{1+\pi_t}\,\frac{M_{t-1}}{P_{t-1}}}_{\text{impuesto inflacionario}} .
  \end{equation}
\item El \concepto{impuesto inflacionario} es la pérdida de valor de los saldos que tiene el público: su base son los saldos reales y su tipo, $\pi/(1+\pi)$. Si los saldos reales no cambian, señoreaje e impuesto coinciden; en tiempo continuo y en estado estacionario, $S=\pi\,M/P$.\footnote{El impuesto inflacionario es antiguo. Los \emph{Económicos} atribuidos a \autor{Aristóteles} (libro II) cuentan que \autor{Dionisio} de Siracusa, endeudado con sus súbditos, mandó entregar las monedas, las reselló con el doble de su valor y pagó con ellas sus deudas.}
\end{itemize}

\begin{anotaciones}
\textbf{Del déficit al impuesto inflacionario, paso a paso} (para consultar; como en el 3.A.39) \vertema{3.A.39}:
\begin{itemize}
\item \emph{Paso 1, el déficit se paga con dinero}: en la restricción consolidada (\ref{sec:politica-fiscal}), si la deuda no cambia, $D_t=D_{t-1}$, todo el déficit se financia con base monetaria nueva:
  \begin{equation*}
  G_t+i_{t-1}D_{t-1}-T_t=M_t-M_{t-1}.
  \end{equation*}
\item \emph{Paso 2, en términos reales}: dividiendo por $P_t$ y multiplicando y dividiendo por $M_t$, con $\gamma_M$ el crecimiento del dinero ($M_{t-1}/M_t=1/(1+\gamma_M)$),
  \begin{equation*}
  \frac{G_t+i_{t-1}D_{t-1}-T_t}{P_t}=\frac{M_t-M_{t-1}}{M_t}\,\frac{M_t}{P_t}=\frac{\gamma_M}{1+\gamma_M}\,\frac{M_t}{P_t}\equiv S_t :
  \end{equation*}
  el señoreaje es el ritmo de emisión por los saldos reales.
\item \emph{Paso 3, con saldos reales constantes}: si la demanda de saldos reales no cambia, $M_t/P_t=M_{t-1}/P_{t-1}$, el dinero y los precios crecen igual, $1+\gamma_M=1+\pi$, y
  \begin{equation*}
  S_t=\frac{\pi}{1+\pi}\,\frac{M_t}{P_t}:
  \end{equation*}
  la monetización funciona como un impuesto con base en los saldos reales y tipo $\pi/(1+\pi)$, el impuesto inflacionario.
\item \emph{Paso 4, en general}: si los saldos reales cambian, el señoreaje se descompone como en~\eqref{eq:senoreaje}: una parte es el aumento de los saldos que el público quiere tener (señoreaje sin inflación, como cuando la economía crece) y otra, el impuesto.
\end{itemize}
\end{anotaciones}

\begin{itemize}
\item La pregunta es por qué la inflación se dispara, si emitir más dinero no da siempre más ingresos. La responde el modelo de \autor{Cagan}.
\end{itemize}

\subsubsection{El modelo de \autor{Cagan} (1956)}\label{sec:cagan}

\paragraph{a) Idea}

\begin{itemize}
\item \textcite{Cagan1956Monetary} quería explicar las hiperinflaciones de entreguerras y de la posguerra con una sola ecuación: la demanda de dinero. Responde a 2 preguntas: ¿puede la monetización continuada de un déficit producir una hiperinflación, aun sin que el dinero se acelere?, y ¿hay un señoreaje máximo? Si el público huye del dinero cuando espera inflación, la base del impuesto inflacionario se encoge, y el Estado tiene que subir el tipo para recaudar lo mismo.
\end{itemize}

\paragraph{b) Modelo}

\subparagraph{i) Supuestos}

\begin{itemize}
\item \emph{La demanda de dinero}: en una hiperinflación la renta y el tipo real cambian poco frente a la inflación, así que la demanda de saldos reales depende sobre todo de la inflación esperada, que es el coste de oportunidad de tener dinero (con la ecuación de \autor{Fisher}, $i\simeq r+\pi^e$):
  \begin{equation}\label{eq:cagan}
  \frac{M^D_t}{P_t}=\zeta\,e^{-\alpha\,\pi^e_t},\qquad\zeta,\alpha>0,
  \end{equation}
  con $\alpha$ la semielasticidad respecto a la inflación esperada y $\zeta$ una constante que recoge la renta y el tipo real. Su función es dar la \emph{base} del impuesto y cómo se encoge.
\item \emph{La oferta de dinero}: la fija la autoridad para financiar el déficit, y crece a la tasa $\gamma_M$. Con el mercado de dinero en equilibrio,
  \begin{equation}\label{eq:cagan-equilibrio}
  \frac{M_t}{P_t}=\zeta\,e^{-\alpha\,\pi^e_t}.
  \end{equation}
  Su función es cerrar el modelo: el nivel de precios se ajusta para que el público quiera tener el dinero emitido.
\item \emph{Las expectativas adaptativas}: en tiempo continuo, la inflación esperada se corrige en proporción al error,
  \begin{equation}\label{eq:cagan-adaptativas}
  \dot\pi^e=\gamma\,\big(\pi-\pi^e\big),\qquad\gamma>0,
  \end{equation}
  con $\gamma$ la velocidad a la que se corrigen (gráfico~\ref{fig:cagan-adaptativas}) \vertema{3.A.5}. Su función es la dinámica: une la inflación de hoy con la esperada de mañana. Cuando \autor{Cagan} escribió, aún no se habían propuesto las expectativas racionales.
\end{itemize}

\begin{figure}[H]
\centering
\caption{Las expectativas adaptativas en el modelo de \autor{Cagan}}\label{fig:cagan-adaptativas}
\grafico[0.6\textwidth]{cagan-adaptativas}
\fuente{Elaboración propia a partir de \textcite{Cagan1956Monetary}. La inflación salta en $t_0$ de $\pi_0$ a $\pi_1$; la esperada la alcanza poco a poco, más deprisa cuanto mayor es $\gamma$.}
\end{figure}

\subparagraph{ii) Desarrollo}

\begin{itemize}
\item \emph{La dinámica de la inflación}. Tomando logaritmos en~\eqref{eq:cagan-equilibrio}, $\ln M_t-\ln P_t=\ln\zeta-\alpha\,\pi^e_t$, y derivando respecto al tiempo (la derivada del logaritmo del dinero es $\gamma_M$, la del de los precios, $\pi$, y $\zeta$ es constante), $\gamma_M-\pi=-\alpha\,\dot\pi^e$: si la inflación esperada sube, el público quiere menos saldos reales y los precios suben más deprisa que el dinero. Con~\eqref{eq:cagan-adaptativas} y despejando:
  \funcionclave[eq:cagan-estabilidad]{\dot\pi^e=\frac{\gamma}{1-\alpha\gamma}\,\big(\gamma_M-\pi^e\big):\quad\text{estable si }\alpha\gamma<1}
  \begin{itemize}
  \item en estado estacionario, $\dot\pi^e=0$, la inflación esperada y la efectiva son iguales al crecimiento del dinero, $\pi=\pi^e=\gamma_M$;
  \item si $\alpha\gamma<1$, las expectativas convergen a $\gamma_M$ y la inflación, al crecimiento del dinero (panel a del gráfico~\ref{fig:cagan-estabilidad}): las hiperinflaciones solo se explican porque el dinero se acelera, que es la conclusión de \autor{Cagan} para las de Europa Central;
  \item si $\alpha\gamma>1$, con una demanda muy sensible y expectativas que se corrigen deprisa, la inflación se aleja de $\gamma_M$ (panel b): la huida del dinero se alimenta a sí misma, sin que el dinero crezca más. Es el refrán del enganche: quien juega con fuego, se quema.
  \end{itemize}
\end{itemize}

\begin{figure}[H]
\centering
\caption{Estabilidad del equilibrio en el modelo de \autor{Cagan} con expectativas adaptativas}\label{fig:cagan-estabilidad}
\grafico[0.85\textwidth]{cagan-estabilidad}
\fuente{Elaboración propia a partir de \textcite{Cagan1956Monetary}. a) $\alpha\gamma<1$: estable. b) $\alpha\gamma>1$: inestable.}
\end{figure}

\begin{anotaciones}
\textbf{La dinámica puede ser inestable: el desarrollo paso a paso} (para consultar; en el cante, el resultado):
\begin{itemize}
\item \emph{Paso 1, el equilibrio en logaritmos}: de~\eqref{eq:cagan-equilibrio}, $\ln M_t-\ln P_t=\ln\zeta-\alpha\,\pi^e_t$.
\item \emph{Paso 2, se deriva respecto al tiempo}:
  \begin{itemize}
  \item $\frac{d\ln M_t}{dt}=\gamma_M$, $\frac{d\ln P_t}{dt}=\pi$ y $\frac{d\ln\zeta}{dt}=0$;
  \item queda $\gamma_M-\pi=-\alpha\,\dot\pi^e$.
  \end{itemize}
\item \emph{Paso 3, se sustituyen las adaptativas}~\eqref{eq:cagan-adaptativas}: $\gamma_M-\pi=-\alpha\gamma\,(\pi-\pi^e)$.
\item \emph{Paso 4, se despeja la inflación efectiva}:
  \begin{itemize}
  \item $\pi\,(1-\alpha\gamma)=\gamma_M-\alpha\gamma\,\pi^e$, es decir, $\pi=\dfrac{\gamma_M-\alpha\gamma\,\pi^e}{1-\alpha\gamma}$;
  \item restando $\pi^e$ a los 2 lados, el error de expectativas es $\pi-\pi^e=\dfrac{\gamma_M-\pi^e}{1-\alpha\gamma}$.
  \end{itemize}
\item \emph{Paso 5, se vuelve a las adaptativas}: $\dot\pi^e=\gamma\,(\pi-\pi^e)=\dfrac{\gamma}{1-\alpha\gamma}\,(\gamma_M-\pi^e)$, que es~\eqref{eq:cagan-estabilidad}.
\item \emph{Paso 6, el estado estacionario}: $\dot\pi^e=0\iff\pi^e=\gamma_M$, y entonces, por el paso 4, $\pi=\gamma_M$.
\item \emph{Paso 7, la estabilidad}: es una ecuación diferencial lineal en $\pi^e$ con pendiente $\frac{d\dot\pi^e}{d\pi^e}=-\frac{\gamma}{1-\alpha\gamma}$:
  \begin{itemize}
  \item si $\alpha\gamma<1$, es negativa: por debajo de $\gamma_M$ la inflación esperada sube, por encima baja, y converge (estable);
  \item si $\alpha\gamma>1$, es positiva: cualquier desviación se agranda (inestable).
  \end{itemize}
\item \emph{Paso 8, la intuición}: si la inflación esperada sube un punto, el público quiere un $\alpha$ por ciento menos de saldos reales, y lo quiere a la velocidad $\gamma$. Si $\alpha\gamma>1$, deshacerse del dinero sube los precios más deprisa de lo que se corrigen las expectativas: la inflación efectiva queda \emph{por encima} de la esperada, las expectativas vuelven a subir y la huida del dinero se acelera.
\end{itemize}
\end{anotaciones}

\begin{itemize}
\item \emph{El señoreaje tiene un máximo}. Hay 2 efectos contrarios: emitir más da más ingresos por unidad de base, pero la inflación encoge la base. Paso a paso:
  \begin{itemize}
  \item \emph{Paso 1, el señoreaje es un tipo por una base}: en tiempo continuo, $S=\dot M/P=(\dot M/M)\,(M/P)=\gamma_M\,M/P$;
  \item \emph{Paso 2, en estado estacionario}: $\pi=\pi^e=\gamma_M$, y la base es la demanda~\eqref{eq:cagan}, $M/P=\zeta e^{-\alpha\pi}$, así que $S(\pi)=\pi\,\zeta\,e^{-\alpha\pi}$;
  \item \emph{Paso 3, la derivada}: $S'(\pi)=\underbrace{\zeta e^{-\alpha\pi}}_{\text{efecto tipo}}\;\underbrace{-\,\alpha\pi\,\zeta e^{-\alpha\pi}}_{\text{efecto base}}=\zeta e^{-\alpha\pi}\,(1-\alpha\pi)$: subir la inflación un punto sube la recaudación en la base que hay, y la baja porque la base cae en una proporción $\alpha\pi$ (la elasticidad de la base a la inflación);
  \item \emph{Paso 4, el máximo}: $S'(\pi)=0\iff\alpha\pi=1$, es decir, $\pi=1/\alpha$; por debajo, $S'>0$ (domina el efecto tipo), y por encima, $S'<0$ (domina el efecto base). La segunda derivada, $S''(\pi)=\alpha\zeta e^{-\alpha\pi}(\alpha\pi-2)$, es negativa en $1/\alpha$: es un máximo;
  \item \emph{Paso 5, el señoreaje máximo}: sustituyendo,
    \funcionclave[eq:laffer-senoreaje]{\pi^{S}=\frac1\alpha,\qquad S^{\max}=S\!\left(\frac1\alpha\right)=\frac{\zeta}{\alpha e}}
    mayor cuanto más dinero se demanda ($\zeta$) y menor cuanto más sensible es la demanda a la inflación ($\alpha$).
  \end{itemize}
  Subir la inflación sube el tipo pero hunde la base: es la \concepto{curva de \autor{Laffer} del señoreaje} (gráfico~\ref{fig:laffer-senoreaje}). Un déficit menor que $S^{\max}$ se puede financiar con 2 inflaciones, $\pi_{\mathrm B}$ y $\pi_{\mathrm A}$; uno mayor, con ninguna. Que $1/\alpha$ maximice el señoreaje no la hace óptima: la inflación tiene los costes del bloque~\ref{sec:efectos}.
\end{itemize}

\begin{figure}[H]
\centering
\caption{La curva de \autor{Laffer} del señoreaje}\label{fig:laffer-senoreaje}
\grafico[0.6\textwidth]{laffer-senoreaje}
\fuente{Elaboración propia a partir de \textcite{Cagan1956Monetary}. El máximo está en $\pi=1/\alpha$; un déficit $S<S^{\max}$ se financia con $\pi_{\mathrm B}$ o con $\pi_{\mathrm A}$, y uno $S'>S^{\max}$, con ninguna.}
\end{figure}

\subparagraph{iii) Implicaciones}

\begin{itemize}
\item La financiación con dinero tiene un techo, $S^{\max}$, y es cara: hacen falta inflaciones altísimas para un señoreaje modesto.
\item Hay 2 caminos a la hiperinflación: un déficit que supera lo que el señoreaje puede dar, que obliga a acelerar el dinero sin fin, o, con $\alpha\gamma>1$, una huida del dinero que alimenta la inflación aunque el dinero no se acelere.
\item Para evitarla hacen falta disciplina fiscal (o un banco central que no financie al Estado) y expectativas ancladas.
\end{itemize}

\paragraph{c) Extensiones}

\begin{itemize}
\item \emph{Expectativas racionales}: el nivel de precios de hoy depende de toda la senda \emph{esperada} del dinero, y un anuncio creíble de que la emisión se acaba baja la inflación ya\nocite{SargentWallace1973Rational}. El desarrollo está en el anexo~\ref{anexo:cagan-racionales}.
\item \emph{El círculo fiscal}: con inflación alta, la recaudación pierde valor por el retraso en el cobro (el efecto \autor{Olivera}-\autor{Tanzi}, \ref{sec:anticipada}), el déficit crece y hay que emitir más; y el público huye hacia la divisa, lo que encoge todavía más la base del impuesto.
\end{itemize}

\paragraph{d) Evidencia empírica}

\begin{itemize}
\item \autor{Cagan} estudió 7 hiperinflaciones (Austria, Alemania, Hungría (2 veces), Polonia, Rusia y Grecia) y concluyó que se explican por el crecimiento continuo de la cantidad de dinero, con una demanda~\eqref{eq:cagan} que se ajustaba bien a los datos.
\item Pero en varias la inflación superó con mucho la que maximizaba el señoreaje, $1/\alpha$: es la \concepto{paradoja de \autor{Cagan}}. La explica la dinámica de las expectativas: mientras el público tarda en corregirlas, acelerar la inflación da más ingresos que mantenerla.
\end{itemize}

\paragraph{e) Valoración}

\begin{itemize}
\item Con una sola ecuación explica el techo del señoreaje y por qué la inflación puede desbocarse. Su punto débil son las expectativas adaptativas: el público se equivoca de forma sistemática durante toda la hiperinflación. Con expectativas racionales, lo que importa es lo que se espera del dinero futuro, y eso explica cómo terminan las hiperinflaciones.
\end{itemize}

\subsubsection{Valoración: el fin de las hiperinflaciones}\label{sec:fin-hiperinflaciones}

\begin{itemize}
\item \textcite{Sargent1982Ends} estudió el fin de 4 hiperinflaciones de los años 20, las de Austria, Hungría, Polonia y Alemania: se pararon de golpe, no poco a poco, cuando cambió el régimen fiscal y monetario (presupuesto equilibrado y un banco central independiente que dejó de financiar al Estado). Los precios se estabilizaron en semanas, con costes de producción pequeños frente a los que predice una curva de \autor{Phillips} con inercia, y la cantidad de dinero siguió creciendo después, porque el público volvió a querer tenerlo.
\item La lección es la de la aritmética monetarista y la teoría fiscal (\ref{sec:politica-fiscal}): la estabilidad de precios exige una política fiscal sostenible y creíble \vertema{3.A.39}. La hiperinflación es fiscal en su origen y se apaga con credibilidad.
\end{itemize}

% --------------------------------------------------------------------------
\subsection{Deflación}\label{sec:deflacion}

\subsubsection{Idea: la deflación buena y la mala}\label{sec:deflacion-idea}

\begin{itemize}
\item La \concepto{deflación} es una caída sostenida y generalizada de los precios. Hay 2 muy distintas\nocite{BordoFilardo2005Deflacion}:
  \begin{itemize}
  \item la \emph{buena}, de oferta: la productividad crece deprisa y los precios bajan con la producción creciendo, como en buena parte del siglo \textsc{xix};
  \item la \emph{mala}, de demanda: el gasto se hunde y los precios bajan con la producción, como en la Gran Depresión.
  \end{itemize}
\item Una caída puntual por el abaratamiento de la energía no es deflación, aunque dé tasas negativas unos meses: en 2020 la inflación de la zona del euro fue negativa de agosto a diciembre y aun así los precios subieron un 0,3~\% en el año\nocite{Eurostat2026IAPC}.
\item \emph{Por qué es peligrosa la deflación de demanda}:
  \begin{itemize}
  \item \emph{deflación por deudas}: las deudas son nominales; si los precios bajan, su valor real sube. Los deudores venden activos y recortan el gasto para pagar, los precios bajan más y la carga real vuelve a subir: «cuanto más pagan los deudores, más deben»\nocite{Fisher1933} \vertema{3.A.42};
  \item \emph{salarios rígidos a la baja}: si los nominales no bajan, los reales suben y el paro aumenta;
  \item \emph{aplazamiento del gasto}: si se espera que los precios bajen, compensa esperar para comprar;
  \item \emph{el límite inferior de los tipos}: el tipo nominal no puede bajar mucho de cero, $i^{\mathrm{ELB}}$, así que con deflación esperada el tipo real, $r=i-\pi^e$, \emph{sube} justo cuando la economía necesita que baje. La deflación hace más probable quedarse en el límite, y el límite hace más probable la deflación \vertema{3.A.37}.
  \end{itemize}
\item \emph{La trampa}: \textcite{Krugman1998Baaack} lo planteó para Japón. Si el tipo natural es negativo y el nominal está en el límite, la única forma de bajar el tipo real es que el público espere inflación; pero un banco central con fama de duro no es creído si la promete, porque al salir de la trampa querrá bajarla: tiene que comprometerse de forma creíble a ser irresponsable \vertema{3.A.40}.\footnote{\autor{Paul Krugman} recibió el Premio Nobel en 2008 «por su análisis de los patrones del comercio y de la localización de la actividad económica»\nocite{Nobel2008Krugman}.} El modelo que lo formaliza es el de \autor{Eggertsson} y \autor{Woodford}.
\end{itemize}

\subsubsection{El modelo de \autor{Eggertsson} y \autor{Woodford} (2003)}\label{sec:eggertsson-woodford}

\paragraph{a) Idea}

\begin{itemize}
\item \textcite{EggertssonWoodford2003Zero} llevan la trampa al modelo canónico de la NEK (\ref{sec:nek-origenes}): ¿qué pasa si una crisis hace negativo el tipo de interés natural y el banco central no puede bajar el nominal por debajo de cero? Explica por qué la deflación de demanda va con recesión y qué puede hacer la política \vertema{3.A.37}.
\end{itemize}

\paragraph{b) Modelo}

\subparagraph{i) Supuestos}

\begin{itemize}
\item La IS dinámica y la NKPC del modelo canónico, sin perturbaciones de costes ($u_t=0$):
  \begin{equation*}
  x_t=\mathbb{E}_t[x_{t+1}]-\frac{1}{\theta}\big(i_t-\mathbb{E}_t[\pi_{t+1}]-r^n_t\big),\qquad\pi_t=\beta\,\mathbb{E}_t[\pi_{t+1}]+\kappa\,x_t .
  \end{equation*}
\item El tipo natural toma 2 valores: en \emph{normalidad}, $r^n>0$; en \emph{crisis}, $r^n_L<0$, por ejemplo tras una crisis financiera en la que los hogares quieren ahorrar y desendeudarse y suben los diferenciales. Si la economía está en crisis, sigue en ella el periodo siguiente con probabilidad $q$ y vuelve a la normalidad con probabilidad $1-q$, y ya no sale de ella ($q$, local).
\item El tipo nominal no puede ser negativo, $i_t\geq0$. En normalidad, el banco central consigue $\pi=x=0$ con $i=r^n>0$, y el límite no es vinculante.
\end{itemize}

\subparagraph{ii) Desarrollo}

\begin{itemize}
\item \emph{Paso 1, en la crisis el límite manda}: haría falta $i=r^n_L<0$; el banco central pone $i=0$. Mientras dura la crisis, la brecha y la inflación son constantes, $x_L$ y $\pi_L$, y como se sigue en crisis con probabilidad $q$ (y en normalidad todo es cero), $\mathbb{E}_t[x_{t+1}]=q\,x_L$ y $\mathbb{E}_t[\pi_{t+1}]=q\,\pi_L$.
\item \emph{Paso 2, se sustituye}:
  \begin{itemize}
  \item de la NKPC, $\pi_L=\beta q\,\pi_L+\kappa\,x_L$, es decir, $\pi_L=\dfrac{\kappa}{1-\beta q}\,x_L$;
  \item de la IS con $i=0$, $(1-q)\,x_L=\dfrac{1}{\theta}\big(q\,\pi_L+r^n_L\big)$.
  \end{itemize}
\item \emph{Paso 3, se despeja}:
  \funcionclave[eq:eggertsson-woodford]{x_L=\frac{(1-\beta q)\,r^n_L}{\theta\,(1-q)(1-\beta q)-q\kappa}<0,\qquad\pi_L=\frac{\kappa\,r^n_L}{\theta\,(1-q)(1-\beta q)-q\kappa}<0}
  con el denominador positivo (la condición para que haya una solución acotada).
\end{itemize}

\subparagraph{iii) Implicaciones}

\begin{itemize}
\item \emph{Recesión y deflación a la vez}, tanto mayores cuanto más negativo es el tipo natural ($r^n_L$) y cuanto más se espera que dure la crisis ($q$, que baja el denominador): la deflación esperada sube el tipo real, $0-q\,\pi_L>0$, que hunde más la demanda, que baja más los precios. Es la deflación por deudas de \autor{Fisher} con la forma de un modelo moderno.
\item \emph{La política monetaria óptima es comprometerse}: prometer que el tipo seguirá en cero \emph{después} de la crisis, con algo de inflación y de brecha positiva, sube hoy la inflación esperada y baja el tipo real. Equivale a un objetivo de nivel de precios, que compensa la deflación pasada \vertema{3.A.36}. Es la base de la orientación sobre los tipos (\emph{forward guidance}) \vertema{3.A.40}.
\item \emph{La política fiscal gana fuerza}: con el tipo en cero, una expansión no sube el tipo y su multiplicador es mayor \vertema{3.A.39}.
\end{itemize}

\paragraph{c) Evidencia empírica}

\begin{itemize}
\item \emph{Japón} vivió una deflación suave pero larga: entre 1998 y 2012 su IPC cayó un 3,9~\% y su deflactor del PIB, un 15,6~\%, con el tipo de interés casi siempre cerca de cero\nocite{BM2026IPCJapon}.
\item La \emph{zona del euro} la rozó: los precios cayeron en tasa interanual entre diciembre de 2014 y marzo de 2015, con un mínimo del $-0{,}6$~\% en enero de 2015, y de nuevo en varios meses de 2016, con una inflación media del 0,2~\% en 2015 y en 2016; en España el IAPC bajó en media un 0,2~\% en 2014, un 0,6~\% en 2015 y un 0,3~\% en 2016\nocite{Eurostat2026IAPC}. El BCE llevó el tipo de la facilidad de depósito a terreno negativo y empezó a comprar deuda pública en 2015 \vertema{3.A.40}.
\item Tras 2008, Estados Unidos y la zona del euro pasaron años en el límite inferior con una inflación baja, pero sin la deflación profunda que predice el modelo.
\end{itemize}

\paragraph{d) Valoración}

\begin{itemize}
\item Explica por qué la deflación de demanda es peligrosa y por qué importan el compromiso y la política fiscal. Pero exagera: predice deflaciones mayores de las observadas y efectos enormes de las promesas sobre tipos muy lejanos (el \emph{enigma de la orientación sobre los tipos}), porque supone expectativas racionales perfectas y precios que miran solo hacia delante.
\item \emph{La salida es reflacionar}, la receta de \autor{Fisher}: devolver los precios a su nivel anterior. Los instrumentos de hoy son la orientación sobre los tipos, las compras de activos \vertema{3.A.40} y la política fiscal: una expansión fiscal que el público no espera pagar con impuestos sube los precios aunque el tipo esté en cero\nocite{Sims2016Fiscal} \vertema{3.A.39}; en el límite, el dinero lanzado desde un helicóptero, que \autor{Ben Bernanke} recordó en 2002\nocite{Bernanke2002Deflation}. Japón lo intentó desde 2013 con un objetivo del 2~\% y compras masivas\nocite{BoJ2013QQE}.
\end{itemize}

\subsubsection{Valoración: desinflar sin caer en deflación}\label{sec:desinflacion}

\begin{itemize}
\item Bajar la inflación tiene un coste: con inercia~\eqref{eq:triangulo-gordon}, la inflación solo baja si el paro sube por encima de su tasa natural. Se mide con la \concepto{ratio de sacrificio}, la producción perdida acumulada, en puntos del PIB, por cada punto que baja la inflación:
  \begin{equation}\label{eq:ratio-sacrificio}
  \text{ratio de sacrificio}=\frac{\sum_t\big(y^n_t-y_t\big)}{\pi_{\text{inicial}}-\pi_{\text{final}}},
  \end{equation}
  con $y$ y $y^n$ la producción y la natural, en logaritmos (por 100).
\item \textcite{Ball1994Sacrifice} la estima en 65 desinflaciones de países de la OCDE: es menor cuanto más rápida es la desinflación y más flexibles son los salarios. La teoría añade que es menor cuanto más creíble es el banco central (\ref{sec:inercia}). Con histéresis, parte del paro se queda y la desinflación sale más cara \vertema{3.A.28}.
\item \emph{Ejemplos}: la Reserva Federal de \autor{Paul Volcker} (presidente de 1979 a 1987) bajó la inflación de Estados Unidos del 14,8~\% de marzo de 1980 al 3,8~\% de diciembre de 1982, con un paro que llegó al 10,8~\% en noviembre de 1982\nocite{BLS2026IPC,BLS2026UNRATE,FedHistory2026Volcker} \vertema{3.A.5};\footnote{\autor{Paul Volcker} (1927-2019) presidió la Reserva Federal de agosto de 1979 a agosto de 1987 y, en 2009-2011, el Consejo Asesor para la Recuperación Económica del presidente Obama\nocite{FedHistory2026Volcker}.} en 2022-2024 la inflación bajó sin recesión, con las expectativas ancladas\nocite{BernankeBlanchard2025What}.
\item Los 2 extremos confirman lo que pedía la parte~\ref{sec:efectos}: una inflación baja, estable y creíble, ni tan alta que el dinero deje de servir ni tan baja que la economía quede atrapada en la deflación. Por eso los bancos centrales prefieren errar algo por arriba.
\end{itemize}

% ==========================================================================
\section*{Conclusión}
\minutos{1}

\begin{itemize}
\item \textbf{\emph{Recapitulación}}:
  \begin{itemize}
  \item En conclusión,
  \begin{itemize}
  \item a largo plazo la inflación es un fenómeno monetario, pero a corto la mueven la demanda, los costes y la inercia, el triángulo de \autor{Gordon}, y la política fiscal y la credibilidad del banco central deciden si el dinero la alimenta;
  \item sus costes son pequeños cuando es baja y anticipada y grandes cuando es alta y volátil; la regla de \autor{Friedman} es el óptimo de un mundo sin rigideces, y las rigideces, el límite inferior y los sesgos de medición llevan a un objetivo algo positivo, el 2~\%;
  \item la hiperinflación nace de un déficit que solo puede pagar la imprenta y termina con una reforma fiscal creíble; la deflación de demanda es peligrosa porque sube la carga de las deudas y ata a la política monetaria.
  \end{itemize}
  \end{itemize}
\item \textbf{\emph{Relevancia}}:
  \begin{itemize}
  \item La inflación de 2021-2023 y el repunte de 2026 recuerdan que una perturbación de costes puede convertirse en inflación persistente si se desanclan las expectativas, y que la inflación tiene costes desiguales.
  \end{itemize}
\item \textbf{\emph{Extensiones}}:
  \begin{itemize}
  \item La demanda de dinero \vertema{3.A.35}, el diseño de la política monetaria y su independencia \vertema{3.A.36}, sus canales y el límite inferior \vertema{3.A.37}, la dominación fiscal \vertema{3.A.39}, la política desde 2008 \vertema{3.A.40}, la nueva economía keynesiana \vertema{3.A.7}, la tasa natural y la histéresis \vertema{3.A.28} y la medición con índices de precios \vertema{3.A.9}.
  \end{itemize}
\item \textbf{\emph{Opinión}}:
  \begin{itemize}
  \item El objetivo simétrico del 2~\% es un compromiso razonable entre los costes de la inflación y los de quedarse sin margen; no se defiende solo con el tipo de interés: lo sostienen la credibilidad del banco central y una política fiscal sostenible.
  \end{itemize}
\item \textbf{\emph{Idea final}}:
  \begin{itemize}
  \item Volviendo al refrán, quien juega con fuego se quema, pero una casa sin fuego se hiela.

  \item El arte de la política monetaria es mantener la lumbre baja y encendida.
  \end{itemize}
\end{itemize}

\clearpage
\appendix

% --------------------------------------------------------------------------
\section{Anexo: los modelos clásico y keynesiano\cortetitulo{} y la historia de la curva de \autor{Phillips}}\label{anexo:clasico-keynesiano}

\subsection{El modelo clásico}\label{anexo:clasico}

\begin{itemize}
\item \emph{Supuestos}: flexibilidad plena de precios, salarios y tipos de interés, que se ajustan hasta vaciar todos los mercados. La economía se describe con 7 ecuaciones\nocite{Serletis2007Demand} \vertema{3.A.4}:
  \begin{enumerate}
  \item demanda de trabajo: el salario real iguala la productividad marginal, $W/P=F_L(K,L)$;
  \item mercado de trabajo en equilibrio, con una oferta que crece con el salario real: $L^D=L^S(W/P)$;
  \item función de producción: $Y=F(K,L)$;
  \item consumo, creciente en la renta disponible y decreciente en el tipo real: $C=C(Y-T,\;i-\pi^e)$;
  \item inversión, decreciente en el tipo real: $I=I(i-\pi^e)$;
  \item equilibrio del mercado de bienes: $Y=C+I+G$;
  \item equilibrio del mercado de dinero, con una demanda de saldos reales creciente en la renta y decreciente en el tipo nominal: $M/P=L(Y,i)$.
  \end{enumerate}
\item \emph{Desarrollo, por bloques}:
  \begin{itemize}
  \item las ecuaciones 1 a 3 fijan el empleo, el salario real y la producción de pleno empleo, $Y_F$, sin que intervenga ninguna variable nominal;
  \item con $Y=Y_F$, las ecuaciones 4 a 6 fijan el tipo de interés real que iguala el ahorro y la inversión;
  \item la 7, con $M$ dado y $i=r+\pi^e$, fija el nivel de precios; con el salario real ya fijado, el salario nominal es $W=(W/P)\,P$.
  \end{itemize}
  En la oferta y la demanda agregadas, la oferta es vertical en $Y_F$ y la demanda solo decide los precios.
\item \emph{Implicaciones}:
  \begin{itemize}
  \item la \concepto{dicotomía clásica}: las variables reales se determinan en las ecuaciones 1 a 6 sin el dinero;
  \item la \concepto{neutralidad}: como $M$ y $P$ solo aparecen juntos, en $M/P$, y $W$ y $P$, en $W/P$, duplicar $M$ duplica $P$ y $W$ y deja todo lo real igual (la diferencia entre las 2 propiedades está en la caja del apartado~\ref{sec:debate-tradicional});
  \item el \concepto{efecto \autor{Fisher}}: con el tipo real fijado por el bloque real, la inflación esperada se traslada uno por uno al tipo nominal, $i=r+\pi^e$ \vertema{3.A.3}.
  \end{itemize}
\item \emph{Las causas de que suban los precios}:
  \begin{itemize}
  \item \emph{de demanda}: una expansión monetaria sube $P$ en la misma proporción, porque es la única forma de que el público quiera tener el dinero nuevo; un aumento del gasto público sube el tipo real y expulsa el mismo consumo e inversión (\emph{crowding-out} total), y como el tipo nominal sube, baja la demanda de saldos reales y sube $P$, pero una sola vez;
  \item \emph{de oferta}: una caída de la productividad o un aumento de la preferencia por el ocio bajan $Y_F$, baja la demanda de saldos reales y sube $P$.
  \end{itemize}
  En todos los casos, la subida desaparece si el impulso no se repite.
\end{itemize}

\subsection{El modelo keynesiano}\label{anexo:keynesiano}

\begin{itemize}
\item \emph{Supuestos}: las mismas ecuaciones, con 3 cambios\nocite{Serletis2007Demand} \vertema{3.A.4} \vertema{3.A.5}:
  \begin{itemize}
  \item el salario nominal no baja de $\overline W$, y el empleo lo fija la demanda de trabajo, $\overline W/P=F_L(K,L)$, con $L^D\leq L^S$: puede haber paro involuntario;
  \item el consumo y la inversión tienen componentes autónomos (el multiplicador y los \emph{animal spirits}): $C=C(C_0,Y-T,i-\pi^e)$ e $I=I(I_0,i-\pi^e)$;
  \item la demanda de dinero es muy sensible al tipo de interés.
  \end{itemize}
\item \emph{Implicaciones}: una expansión monetaria sube $P$, baja el salario real $\overline W/P$ y sube el empleo y la producción: se rompen la dicotomía y la neutralidad. Una expansión fiscal sube la producción y el tipo, con un \emph{crowding-out} parcial, y puede subir algo los precios. La oferta agregada tiene pendiente positiva: los precios suben menos que en el modelo clásico, porque parte del impulso va a la producción. Con rigideces nominales y reales, el efecto sobre los precios persiste aunque se retire el estímulo.
\end{itemize}

\begin{anotaciones}
\textbf{Sobre la fuente}: el desarrollo de estos modelos sigue los capítulos 1 («Classical Macroeconomic Theory»: el modelo clásico completo, la dicotomía, la oferta y la demanda agregadas y la neutralidad del dinero) y 2 de \textcite{Serletis2007Demand}.
\end{anotaciones}

\subsection{La curva de \autor{Phillips} antes de \autor{Friedman}}\label{anexo:phillips-historia}

\begin{itemize}
\item \textcite{Phillips1958}, con datos del Reino Unido de 1861-1913, estimó \vertema{3.A.28}
  \begin{equation*}
  \log_{10}(\gamma_W+0{,}900)=0{,}984-1{,}394\,\log_{10}U\quad\Longleftrightarrow\quad\gamma_W=9{,}638\,U^{-1{,}394}-0{,}900,
  \end{equation*}
  con $\gamma_W$ el crecimiento anual del salario nominal y $U$ la tasa de paro, en~\%. Con un paro de alrededor del 5,5~\% los salarios no crecen; con algo menos del 2,5~\% (el 2,4~\% según la ecuación) crecen un 2~\%, lo que crecía la productividad, y los precios quedan estables. Los datos de 1948-1957 encajaban en la curva de 1861-1913.
\item \textcite{Lipsey1960Relation} le dio fundamento: el salario sube cuando la demanda de trabajo supera a la oferta, $\gamma_W=g[(L^D-L^S)/L^S]$, $g'>0$, y el exceso de demanda baja con el paro, $(L^D-L^S)/L^S=h(U)$, $h'<0$; juntando, $\gamma_W=g[h(U)]\equiv f(U)$, $f'<0$. Con exceso de demanda nulo queda el paro friccional. Añadió a la regresión la variación del paro y la de los precios al consumo.
\item \textcite{SamuelsonSolow1960Analytical} la pasaron a precios, con un margen constante sobre los costes laborales unitarios ($\pi=\gamma_W-\gamma_{\mathrm{PAT}}$), y la leyeron como un \emph{menú} entre inflación y paro \vertema{3.A.5}.
\item \emph{Valoración}: no tenía expectativas ni explicaba la inflación de oferta, y el menú resultó inestable en los 70. Pero desde entonces todos los modelos han tenido que explicar la relación\nocite{Akerlof2002Behavioral}.
\end{itemize}

% --------------------------------------------------------------------------
\section{Anexo: información imperfecta, inconsistencia temporal\cortetitulo{} y política fiscal}\label{anexo:otras-explicaciones}

\subsection{El modelo de islas de \autor{Lucas}}\label{anexo:lucas}

\begin{itemize}
\item \emph{El problema de cada productor}: ve el precio de su isla, $p_{i,t}=p_t+z_{i,t}$, con $z_{i,t}$ una perturbación real de media cero y varianza $\sigma^2_z$, independiente del nivel general $p_t$, que tiene media $p^e_t$ y varianza $\sigma^2_p$\nocite{Lucas1972Expectations,Lucas1973Some} \vertema{3.A.6}.
\item \emph{La extracción de la señal}: con expectativas racionales (y normalidad), su mejor estimación de la perturbación real es
  \begin{equation*}
  \mathbb{E}[z_{i,t}\mid p_{i,t}]=\omega\,\big(p_{i,t}-p^e_t\big),\qquad\omega=\frac{\sigma^2_z}{\sigma^2_z+\sigma^2_p}:
  \end{equation*}
  atribuye a causas reales la parte $\omega$ de la diferencia entre su precio y el nivel esperado, el \emph{coeficiente de proyección}.
\item \emph{La oferta de cada isla}: produce más cuanto mayor cree que es su precio relativo, $y_{i,t}=y^n+\varepsilon_S\,\omega\,(p_{i,t}-p^e_t)$.
\item \emph{Agregando}: como las $z_{i,t}$ suman cero, la oferta de \autor{Lucas}~\eqref{eq:oferta-lucas} es $y_t=y^n+\psi\,(p_t-p^e_t)$, con $\psi\equiv\omega\,\varepsilon_S$. Solo la parte no esperada de los precios mueve la producción, y menos cuanto mayor es la varianza de los precios: con inflación errática, $\omega$ y $\psi$ se acercan a cero y la curva de \autor{Phillips} se vuelve vertical.
\item \emph{Con la demanda} $y_t=m_t-p_t$, solo la parte no anticipada del dinero tiene efectos reales, la proposición de inefectividad\nocite{SargentWallace1975Rational}, y como $\psi$ cambia con la política, la forma reducida no sirve para evaluarla (la crítica de \autor{Lucas}).\findemo
\end{itemize}

\subsection{\autor{Kydland} y \autor{Prescott} y \autor{Barro} y \autor{Gordon}}\label{anexo:barro-gordon}

\begin{itemize}
\item \emph{Supuestos}\nocite{KydlandPrescott1977Rules,BarroGordon1983Positive} \vertema{3.A.36}: la curva de \autor{Phillips} con sorpresas, $U_t=U^n-\alpha\,(\pi_t-\pi^e_t)$, con $\alpha=1/a$; un banco central que minimiza la pérdida
  \begin{equation*}
  L_t=(U_t-U^*)^2+\gamma\,(\pi_t-\pi^*)^2,\qquad U^*<U^n,
  \end{equation*}
  y un juego de información completa en 2 etapas: primero el público fija $\pi^e_t$ y después el banco central elige $\pi_t$.
\item \emph{Discreción}: dada $\pi^e_t$, sustituyendo la curva en la pérdida y derivando respecto a $\pi_t$, $-2\alpha\,(U_t-U^*)+2\gamma\,(\pi_t-\pi^*)=0$: la condición de primer orden es \emph{lineal}, $\gamma\,(\pi_t-\pi^*)=\alpha\,(U_t-U^*)$. Con expectativas racionales, $\pi^e_t=\pi_t$, el paro queda en $U^n$ y la inflación en $\pi^{\mathrm D}=\pi^*+\frac{\alpha}{\gamma}(U^n-U^*)$, la ecuación~\eqref{eq:sesgo}.
\item \emph{Los 3 puntos}: la \emph{regla} (R: $\pi^*$, $U^n$), la \emph{tentación} (T: anunciada la regla y creída, el banco central sorprende, $\pi^T=\pi^*+\alpha(U^n-U^*)/(\gamma+\alpha^2)$ y $U<U^n$) y la \emph{discreción} (D). Las pérdidas cumplen $L^T<L^R<L^D$: la regla no es creíble porque la tentación es mejor para el banco central, y el público, que lo sabe, lleva la economía a la discreción, la peor de las 3.
\item \emph{Soluciones}: reglas con compromiso; reputación en el juego repetido \vertema{3.A.14}; un banco central conservador, con un $\gamma$ mayor que el de la sociedad\nocite{Rogoff1985Optimal}; atarse las manos con un tipo de cambio fijo frente a un país de baja inflación\nocite{GiavazziPagano1988Advantage}; independencia y objetivos de inflación. Que la pérdida del banco central se pueda derivar de la utilidad de los hogares lo mostraron después \textcite{RotembergWoodford1997Optimization} en la nueva economía keynesiana.
\end{itemize}

\subsection{La aritmética monetarista y la teoría fiscal}\label{anexo:fiscal}

\begin{itemize}
\item Se parte de la restricción consolidada en porcentaje del PIB (caja del apartado~\ref{sec:politica-fiscal}), sin inflación no esperada \vertema{3.A.39}:
  \begin{equation*}
  d_t-d_{t-1}=(g_t-\tau_t)-s_t+\frac{r-\gamma_Y}{1+\gamma_Y}\,d_{t-1}.
  \end{equation*}
\item \emph{La aritmética monetarista desagradable}\nocite{SargentWallace1981}: si la deuda no puede superar un nivel desde la fecha $\bar t$ y la política fiscal no se ajusta, el señoreaje necesario desde entonces, $s^{\text{después}}$, supera al que habría bastado desde el principio, $s^*\equiv(g-\tau)+\frac{r-\gamma_Y}{1+\gamma_Y}\,d_0$:
  \begin{equation*}
  s^{\text{después}}-s^*=\left[\left(\frac{1+r}{1+\gamma_Y}\right)^{\bar t}-1\right]\left(s^*-s^{\text{antes}}\right)>0 ,
  \end{equation*}
  con $s^{\text{antes}}<s^*$ el señoreaje del periodo de dureza monetaria y $r>\gamma_Y$. Cada punto de señoreaje que se ahorra hoy cuesta más señoreaje mañana.
\item \emph{La teoría fiscal}: con superávits futuros dados, \eqref{eq:valoracion} deja de ser una restricción y fija el nivel de precios; en el régimen de \textcite{Leeper1991}, una política monetaria pasiva y una fiscal activa.
\end{itemize}

% --------------------------------------------------------------------------
\section{Anexo: las rigideces de la nueva economía keynesiana,\cortetitulo{} la curva de \autor{Phillips} y la persistencia}\label{anexo:persistencia}

\subsection{Cómo modela la NEK las rigideces: \autor{Fischer}, \autor{Taylor} y \autor{Calvo}}\label{anexo:rigideces}

\begin{anotaciones}
\textbf{De \autor{Fischer} y \autor{Taylor} a \autor{Calvo}: por qué la 2.ª generación se queda con \autor{Calvo}} (para consultar; no se canta; como en el 3.A.7) \vertema{3.A.7}:
\begin{itemize}
\item \emph{Para qué nacen los contratos escalonados}: los 2 modelos responden a la proposición de inefectividad de \textcite{SargentWallace1975Rational}: con expectativas racionales, la política monetaria sistemática no tendría efectos reales \vertema{3.A.6}.
\item \emph{\autor{Fischer} (1977)}\nocite{Fischer1977Long}:
  \begin{itemize}
  \item los contratos duran 2 periodos y se firman por mitades;
  \item el salario de cada periodo se fija al firmar, con lo que se espera entonces: $w^j_t=\mathbb{E}_{t-j}[p_t]$, y puede ser distinto en cada periodo del contrato;
  \item el banco central puede reaccionar a lo que pasa después de firmados los contratos y estabilizar la producción, aunque las expectativas sean racionales.
  \end{itemize}
\item \emph{\autor{Taylor} (1980)}\nocite{Taylor1980Aggregate}:
  \begin{itemize}
  \item el salario queda fijo durante todo el contrato y las revisiones no están sincronizadas;
  \item cada empresa que revisa mira el salario de las que no revisan;
  \item los efectos de la política monetaria duran más que los propios contratos, aunque al final se ajustan los precios.
  \end{itemize}
\item \emph{Lo que muestran los 2}: expectativas racionales y política eficaz son compatibles. Lo que da la neutralidad es la flexibilidad de los precios, no las expectativas.
\item \emph{Por qué los superó la regla de \textcite{Calvo1983Staggered}} y encaja mejor en los modelos de equilibrio general dinámicos y estocásticos:
  \begin{itemize}
  \item \emph{no tiene memoria}: cada empresa revisa con la misma probabilidad, $1-\eta$, haga el tiempo que haga, así que no hace falta seguir a cada cohorte de contratos: todo lo que importa del pasado está en el nivel de precios del periodo anterior, $\pi_t=(1-\eta)(p^*_t-p_{t-1})$ (gráfico~\ref{fig:reglas-precios});
  \item \emph{sale de optimizar}: la empresa que revisa elige el precio que maximiza el valor esperado de sus beneficios mientras dure; las reglas de \autor{Fischer} y de \autor{Taylor} se suponían;
  \item \emph{se agrega en una ecuación}: la curva de \autor{Phillips} neokeynesiana, con una pendiente que depende de parámetros estructurales;
  \item \emph{da el coste de la inflación}: la inflación dispersa los precios relativos, y esa dispersión es la que justifica que el banco central se preocupe por la inflación (anexo~\ref{anexo:inflacion-optima});
  \item \emph{su precio}: la frecuencia de los cambios es exógena y no sube con la inflación, como sí predicen las reglas dependientes del estado.
  \end{itemize}
\end{itemize}
\begin{figuracaja}
\caption{Cómo se fijan los precios en cada modelo ante una subida permanente del dinero}\label{fig:reglas-precios}
\grafico[0.95\textwidth]{reglas-precios}
\fuente{Elaboración propia a partir de \textcite{Fischer1977Long}, \textcite{Taylor1980Aggregate}, \textcite{Calvo1983Staggered} y \autor{Rotemberg} (1982); trayectorias ilustrativas, no calibradas.}
\end{figuracaja}
\end{anotaciones}

\subsection{De la rigidez de los precios a la curva de \autor{Phillips}}\label{anexo:nkpc}

\begin{itemize}
\item \emph{La regla de \autor{Calvo}}: cada periodo, cada empresa mantiene su precio con probabilidad $\eta$ (duración media $1/(1-\eta)$) y lo revisa con probabilidad $1-\eta$.
\item \emph{Agregación} (en logaritmos): el nivel de precios mezcla los que revisan, que eligen $p^*_t$, y los que no, cuyo precio medio es el del periodo anterior, $p_t=\eta\,p_{t-1}+(1-\eta)\,p^*_t$; es decir, $\pi_t=(1-\eta)\,(p^*_t-p_{t-1})$.
\item \emph{El precio elegido}: como puede quedarse fijo, es un margen sobre la media descontada de los costes marginales nominales que se esperan mientras dure, $p^*_t=\mu+(1-\beta\eta)\sum_{k\geq0}(\beta\eta)^k\,\mathbb{E}_t[\mathrm{cmg}^n_{t+k}]$, con $\mu$ el margen en logaritmos.
\item Juntando las 2\nocite{Gali2015Monetary}, $\pi_t=\beta\,\mathbb{E}_t[\pi_{t+1}]+\frac{(1-\eta)(1-\beta\eta)}{\eta}\,\widehat{\mathrm{cmg}}_t$, y con el coste marginal real proporcional a la brecha, la NKPC~\eqref{eq:nkpc} con la pendiente $\kappa\equiv\frac{(1-\eta)(1-\beta\eta)}{\eta}\,(\theta+\varphi)$, con $\varphi$ la inversa de la elasticidad de la oferta de trabajo: cuanto más rígidos los precios ($\eta$ mayor), más plana \vertema{3.A.7}.
\item \emph{La NAIRU en la NKPC}: con la ley de \autor{Okun}, la brecha es proporcional a $-(U_t-U^n)$ y $\pi_t=\beta\,\mathbb{E}_t[\pi_{t+1}]-a\,(U_t-U^n)+u_t$. Con el paro por debajo de $U^n$, $\pi_t>\beta\,\mathbb{E}_t[\pi_{t+1}]$: la inflación es alta y se espera que \emph{baje}\nocite{Mankiw2001Inexorable}. El nombre de NAIRU viene de la curva aceleracionista~\eqref{eq:aceleracionista}, no de esta (\ref{sec:inercia}).
\end{itemize}

\subsection{Modelos con persistencia intrínseca}\label{anexo:persistencia-modelos}

\begin{itemize}
\item \emph{Indexación}\nocite{ChristianoEichenbaumEvans2005}: si los precios que no se revisan se actualizan con la inflación del periodo anterior, la NKPC se cumple para la inflación \emph{en exceso} de la pasada,
  \begin{equation*}
  \pi_t-\pi_{t-1}=\beta\,\big(\mathbb{E}_t[\pi_{t+1}]-\pi_t\big)+\kappa\,x_t+u_t ,
  \end{equation*}
  y despejando $\pi_t$,
  \begin{equation*}
  \pi_t=\frac{1}{1+\beta}\,\pi_{t-1}+\frac{\beta}{1+\beta}\,\mathbb{E}_t[\pi_{t+1}]+\frac{\kappa}{1+\beta}\,x_t+\frac{u_t}{1+\beta}:
  \end{equation*}
  es la híbrida~\eqref{eq:nkpc-hibrida} con $\gamma_b=1/(1+\beta)$ y $\gamma_f=\beta/(1+\beta)$ (en el cuerpo se escribe $\kappa$ por sencillez para la pendiente dividida por $1+\beta$).\findemo
\item \emph{Reglas sencillas}\nocite{GaliGertler1999Inflation}: una parte de las empresas que revisan fija el precio con una regla que mira la inflación pasada; sale también una curva híbrida, en la que el peso retrospectivo es la proporción de esas empresas, y la estiman con el coste marginal real en lugar de la brecha.
\item \emph{Contratos sobre salarios relativos}\nocite{FuhrerMoore1995Inflation}: cada trabajador negocia un salario de contrato que mira a los salarios reales de los contratos vigentes, los anteriores y los que vendrán; con contratos de 2 periodos, la inflación depende con pesos iguales de la pasada y de la esperada, más la brecha, y hereda la inercia de los salarios.
\item \emph{Información rígida}\nocite{MankiwReis2002Sticky}: cada periodo una fracción $\lambda$ de las empresas pone al día su información; todas cambian el precio, pero las demás lo fijan con lo que sabían cuando se informaron. La curva resultante es
  \begin{equation*}
  \pi_t=\frac{\xi\lambda}{1-\lambda}\,x_t+\lambda\sum_{j\geq0}(1-\lambda)^j\,\mathbb{E}_{t-1-j}\big[\pi_t+\xi\,\Delta x_t\big],
  \end{equation*}
  con $\lambda$ y $\xi$ (la sensibilidad del precio deseado a la brecha) locales. Lo que pesa son las expectativas formadas \emph{en el pasado} sobre el presente: las perturbaciones pasan a los precios poco a poco y una desinflación anunciada cuesta producción.
\end{itemize}

\subsection{Otras fuentes de inercia}\label{anexo:inercia-otras}

\begin{itemize}
\item Las expectativas: si son adaptativas, o si el público duda del objetivo, la inflación esperada sigue a la pasada.
\item La política económica: el banco central tarda en reaccionar y sus decisiones tardan en hacer efecto.
\item Con mucha indexación y pocas empresas que fijen el precio de forma óptima, estabilizar la inflación exige movimientos muy grandes de los tipos de interés a lo largo del ciclo.
\end{itemize}

% --------------------------------------------------------------------------
\section{Anexo: la inflación óptima con impuestos distorsionantes\cortetitulo{} y con precios rígidos}\label{anexo:inflacion-optima}

\subsection{Impuestos distorsionantes: la inflación como impuesto}\label{anexo:optima-impuestos}

\begin{itemize}
\item \emph{Sin impuestos de suma fija}: si los demás impuestos distorsionan, la inflación es un impuesto más, sobre los saldos reales. La imposición óptima (\autor{Ramsey}) iguala el coste marginal de recaudar con cada impuesto, y pide usar un poco todos, también la inflación: $\pi^*>-r$\nocite{Phelps1973Inflation}. Es la pregunta de cuánto usar el impuesto inflacionario \vertema{3.A.39}.
\item \emph{Pero la regla resiste}: si el dinero sirve para comprar (es un bien intermedio, como en los modelos de pago por adelantado) y el consumo ya se grava, gravar el dinero es gravar 2 veces lo mismo; como los bienes intermedios no deben gravarse, la regla de \autor{Friedman} vuelve a ser óptima aunque los demás impuestos distorsionen\nocite{ChariChristianoKehoe1996Optimality}.
\item \emph{Rentas que escapan al fisco}: si hay beneficios puros o economía sumergida que no se pueden gravar, la inflación es un modo indirecto de alcanzarlos; pero para justificar un 2~\% la renta no gravada tendría que superar el 30~\% del total\nocite{SchmittGroheUribe2010Optimal}.
\item \emph{Demanda exterior de la moneda}: un país cuya moneda se usa fuera cobra con la inflación a los extranjeros; para Estados Unidos saldría una inflación óptima del 2 al 10~\%. Pero países cuya moneda casi nadie usa fuera tienen objetivos parecidos\nocite{SchmittGroheUribe2010Optimal}.
\end{itemize}

\subsection{Precios rígidos: la estabilidad de precios}\label{anexo:optima-rigideces}

\begin{itemize}
\item \emph{La dispersión de los precios}: con precios à la \autor{Calvo}, en cada momento conviven precios fijados en fechas distintas. Con inflación (positiva o negativa), los que no se han revisado se quedan atrás o por delante, y bienes que los hogares valoran igual tienen precios distintos: compran cantidades distintas y se produce de forma ineficiente.
\item \emph{La pérdida social}: una aproximación de segundo orden a la utilidad de los hogares da una pérdida cuadrática en la inflación y en la brecha\nocite{Woodford2003Interest,Gali2015Monetary} \vertema{3.A.7}:
  \begin{equation*}
  L=\mathbb{E}_0\sum_{t=0}^{\infty}\beta^t\big(\pi_t^2+\lambda_x\,x_t^2\big),\qquad\lambda_x=\frac{\kappa}{\sigma},
  \end{equation*}
  con $\sigma$ la elasticidad de sustitución entre variedades. La inflación pesa mucho porque cuanto más sustituibles son las variedades ($\sigma$ alta), más cuesta la dispersión.
\item \emph{La inflación óptima es cero}: cualquier inflación media distinta de cero mantiene una dispersión permanente, así que el óptimo es la \concepto{estabilidad de precios}, $\pi^*=0$, y no la regla de \autor{Friedman}.
\item \emph{Con las 2 fricciones}, la del dinero (que pide $-r$) y la de los precios (que pide 0), el óptimo queda entre las 2 y, con parámetros razonables, muy cerca de cero: en un modelo estimado para Estados Unidos, una media del $-0{,}4$~\%\nocite{SchmittGroheUribe2010Optimal}. Las rigideces nominales pesan mucho más que el coste de tener dinero.\findemo
\end{itemize}

% --------------------------------------------------------------------------
\section{Anexo: el modelo de \autor{Cagan} con expectativas racionales}\label{anexo:cagan-racionales}

\begin{itemize}
\item \emph{Paso 1, la demanda en tiempo discreto}: en logaritmos ($m_t$ dinero y $p_t$ nivel de precios; aquí $m$ es un logaritmo, no los saldos reales) y con la inflación esperada $\mathbb{E}_t[p_{t+1}]-p_t$, la demanda~\eqref{eq:cagan} en equilibrio es\nocite{SargentWallace1973Rational}
  \begin{equation*}
  m_t-p_t=\ln\zeta-\alpha\,\big(\mathbb{E}_t[p_{t+1}]-p_t\big).
  \end{equation*}
\item \emph{Paso 2, se despeja el nivel de precios}: pasando $p_t$ a un lado, $p_t\,(1+\alpha)=m_t-\ln\zeta+\alpha\,\mathbb{E}_t[p_{t+1}]$, es decir,
  \begin{equation}\label{eq:cagan-precios}
  p_t=\frac{1}{1+\alpha}\,\big(m_t-\ln\zeta\big)+\frac{\alpha}{1+\alpha}\,\mathbb{E}_t[p_{t+1}].
  \end{equation}
  El nivel de precios de hoy depende de 2 cosas: el dinero de hoy y el nivel de precios que se espera mañana. No es aún una solución, porque a la derecha hay una variable endógena, la expectativa: falta decir cómo se forma.
\end{itemize}

\subsection{Con expectativas adaptativas}\label{anexo:cagan-hea}

\begin{itemize}
\item \emph{Paso 3}: la inflación esperada es una media ponderada de la pasada, $\mathbb{E}_t[p_{t+1}]-p_t=\gamma\sum_{j\geq0}(1-\gamma)^j\pi_{t-j}$, con $0<\gamma\leq1$ la capacidad de aprender de los errores.
\item \emph{Paso 4}: llevándola a la demanda del paso 1,
  \begin{equation*}
  p_t=m_t-\ln\zeta+\alpha\gamma\sum_{j\geq0}(1-\gamma)^j\,\pi_{t-j}.
  \end{equation*}
\item \emph{Implicación}: el nivel de precios depende del dinero de hoy y de toda la historia de la inflación; el futuro no cuenta, y un anuncio sobre el dinero futuro no cambia nada hoy («hasta que no lo veo no lo creo»).
\end{itemize}

\subsection{Con expectativas racionales}\label{anexo:cagan-her}

\begin{itemize}
\item \emph{Paso 5, la ecuación un periodo después}: escribiendo~\eqref{eq:cagan-precios} en $t+1$ y tomando la esperanza en $t$, con la ley de las expectativas iteradas, $\mathbb{E}_t\big[\mathbb{E}_{t+1}[p_{t+2}]\big]=\mathbb{E}_t[p_{t+2}]$:
  \begin{equation*}
  \mathbb{E}_t[p_{t+1}]=\frac{1}{1+\alpha}\,\big(\mathbb{E}_t[m_{t+1}]-\ln\zeta\big)+\frac{\alpha}{1+\alpha}\,\mathbb{E}_t[p_{t+2}].
  \end{equation*}
\item \emph{Paso 6, se sustituye en~\eqref{eq:cagan-precios}}:
  \begin{equation*}
  p_t=\frac{1}{1+\alpha}\left[(m_t-\ln\zeta)+\frac{\alpha}{1+\alpha}\,\big(\mathbb{E}_t[m_{t+1}]-\ln\zeta\big)\right]+\left(\frac{\alpha}{1+\alpha}\right)^{2}\mathbb{E}_t[p_{t+2}],
  \end{equation*}
  que aún tiene la expectativa del precio en $t+2$.
\item \emph{Paso 7, se itera $J$ veces}:
  \begin{equation*}
  p_t=\frac{1}{1+\alpha}\sum_{j=0}^{J-1}\left(\frac{\alpha}{1+\alpha}\right)^{j}\big(\mathbb{E}_t[m_{t+j}]-\ln\zeta\big)+\left(\frac{\alpha}{1+\alpha}\right)^{J}\mathbb{E}_t[p_{t+J}].
  \end{equation*}
\item \emph{Paso 8, sin burbujas}: si el último término tiende a cero, $\lim_{J\to\infty}\left(\alpha/(1+\alpha)\right)^J\mathbb{E}_t[p_{t+J}]=0$, y como los pesos $\frac{1}{1+\alpha}\left(\frac{\alpha}{1+\alpha}\right)^j$ suman 1,
  \begin{equation}\label{eq:cagan-racionales}
  p_t=\frac{1}{1+\alpha}\sum_{j\geq0}\left(\frac{\alpha}{1+\alpha}\right)^{j}\mathbb{E}_t[m_{t+j}]-\ln\zeta .
  \end{equation}
\item \emph{Paso 9, con el dinero creciendo a tasa constante}: $m_{t+j}=m_t+j\,\gamma_M$, y como $\frac{1}{1+\alpha}\sum_{j\geq0}\left(\frac{\alpha}{1+\alpha}\right)^j j=\alpha$, \eqref{eq:cagan-racionales} da $p_t=m_t+\alpha\gamma_M-\ln\zeta$: la inflación es $\gamma_M$ y los saldos reales, $\ln\zeta-\alpha\gamma_M$ en logaritmos, los de la curva de \autor{Laffer}~\eqref{eq:laffer-senoreaje}.
\item \emph{Paso 10, un plan de estabilización creíble}: si en $t$ el Gobierno anuncia, y el público lo cree, que desde mañana el dinero deja de crecer ($m_{t+j}=m_t$), \eqref{eq:cagan-racionales} da $p_t=m_t-\ln\zeta$: la inflación cae a cero \emph{el día del anuncio}, y los saldos reales que el público quiere suben de $\ln\zeta-\alpha\gamma_M$ a $\ln\zeta$. Para que los precios no bajen, el banco central puede aumentar el dinero una sola vez en $\alpha\gamma_M$ sin causar inflación: es la \emph{remonetización} que siguió al fin de las hiperinflaciones.
\item \emph{Implicaciones}:
  \begin{enumerate}
  \item el nivel de precios es una media ponderada de toda la senda \emph{esperada} del dinero, con pesos que bajan con el horizonte: el pasado no importa, el futuro sí;
  \item lo que mueve los precios es lo que el público \emph{cree} sobre el dinero futuro, no la senda que de verdad siga, ni lo que se anuncie: la credibilidad es decisiva. Si un plan de estabilización es creíble, los precios dejan de subir el día del anuncio y los saldos reales se recuperan aunque el dinero siga creciendo un tiempo, que es lo que \textcite{Sargent1982Ends} encontró en los años 20; si no lo es, no consigue nada;
  \item sin la condición de no burbujas caben otras soluciones, $p_t$ más una burbuja $b_t$ con $\mathbb{E}_t[b_{t+1}]=\frac{1+\alpha}{\alpha}\,b_t$, que crece sin límite: sendas que se alejan del fundamento, otra forma de la inestabilidad de~\eqref{eq:cagan-estabilidad}.\findemo
  \end{enumerate}
\end{itemize}

% --------------------------------------------------------------------------
\section{Anexo: la inflación y las cuentas públicas}\label{anexo:cuentas-publicas}

\begin{itemize}
\item A corto plazo la inflación mejora las cuentas públicas: la recaudación sube con los precios (y más si los impuestos no se indexan, \ref{sec:anticipada}), muchos gastos se actualizan con retraso y la inflación no esperada reduce el valor real de la deuda emitida a tipo fijo \vertema{3.A.39}.
\item Con el tiempo el efecto se va (gráfico~\ref{fig:inflacion-cuentas-publicas}): se indexan las pensiones y los salarios públicos, la deuda se renueva a tipos más altos, el banco central sube los tipos y el Gobierno aprueba medidas para compensar a los hogares\nocite{FMI2023FiscalMonitor}.
\end{itemize}

\begin{figure}[H]
\centering
\caption{Efectos de la inflación sobre las cuentas públicas}\label{fig:inflacion-cuentas-publicas}
\grafico[0.8\textwidth]{inflacion-cuentas-publicas}
\fuente{Elaboración propia a partir de \textcite{FMI2023FiscalMonitor}.}
\end{figure}

% --------------------------------------------------------------------------
\section{Anexo: preguntas de test}\label{anexo:test}

\begin{anotaciones}
Preguntas tipo test del primer ejercicio: \href{https://www.victorgutierrezmarcos.es/oposicion/temario/primer-ejercicio/test/simulador.html?tema=3.A.41}{simulador de test}.
\end{anotaciones}

\bibliografia
\end{document}
